% main.tex — Fee Concentration Insurance Specification
% Master document: \input{} for all components

\documentclass[11pt,a4paper]{article}

% preamble.tex — Shared packages, notation, theorem environments
% Fee Concentration Insurance Specification

% ── Packages ──────────────────────────────────────────────────────────────────
\usepackage{amsmath,amssymb,amsthm}
\usepackage{mathtools}
\usepackage{hyperref}
\usepackage{cleveref}
\usepackage{booktabs}
\usepackage{longtable}
\usepackage{enumitem}

% ── Theorem environments ──────────────────────────────────────────────────────
\newtheorem{theorem}{Theorem}[section]
\newtheorem{proposition}[theorem]{Proposition}
\newtheorem{lemma}[theorem]{Lemma}
\newtheorem{corollary}[theorem]{Corollary}
\newtheorem{definition}[theorem]{Definition}
\newtheorem{remark}[theorem]{Remark}

% ── Invariant environment (Hoare triples) ─────────────────────────────────────
\newcounter{invariant}
\newenvironment{invariant}[2]{% #1 = ID (e.g., CFMM-01), #2 = Name
  \refstepcounter{invariant}%
  \par\medskip\noindent\textbf{#1} (\textit{#2}).\quad
}{\par\medskip}

% ── Fee concentration index (primary state) ──────────────────────────────────
\newcommand{\AT}{A_T}                       % fee concentration index (primary, tracked)
\newcommand{\ThetaSum}{\Theta}              % theta sum: Σ θ_k (co-primary, tracked)
\newcommand{\PosCount}{N}                   % position count (co-primary, tracked)
\newcommand{\ATnull}{A_T^{1/\PosCount}}     % Ma-Crapis null: sqrt(Θ/N²)
\newcommand{\DeltaPlus}{\Delta^{+}}         % concentration deviation: max(0, A_T - A_T^{1/N})
\newcommand{\DeltaStar}{\Delta^{*}}         % econometric turning point (≈ 0.09)
\newcommand{\thetapos}{\theta}              % position sophistication weight
\newcommand{\blocklife}{\ell}               % position lifetime in blocks

% ── Pricing ──────────────────────────────────────────────────────────────────
% p = Δ⁺/(1-Δ⁺)    odds-ratio from deviation to price space
% p ∈ [0, ∞)        for Δ⁺ ∈ [0, 1)

% ── Fees and funding ─────────────────────────────────────────────────────────
\newcommand{\fee}{\varphi}                  % trading fee φ
\newcommand{\feebase}{\varphi_{\text{base}}}
\newcommand{\feefund}{\varphi_{\text{fund}}}
\newcommand{\fundrate}{r}                   % funding rate

% ── Insurance CFMM notation (power² payoff) ─────────────────────────────────
\newcommand{\premfactor}{\gamma}            % base premium factor (borrow rate floor)
\newcommand{\Res}{R}                         % per-pool reserve balance
\newcommand{\protgrowth}{g^{\text{prot}}}   % protection growth per liquidity (Q128)
\newcommand{\donateamt}{D}                   % donate amount at a liquidity event

% ── Trading function (linearized power²) ────────────────────────────────────
% \psi is a standard LaTeX command — use it directly for the trading function ψ(u,y)
\newcommand{\pl}{p_l}                        % lower tick price (effective strike)
\newcommand{\pu}{p_u}                        % upper tick price (cap)
\newcommand{\pstar}{p^{*}}                   % strike price = Δ*/(1-Δ*)
\newcommand{\pmark}{p_{\text{mark}}}         % CFMM mark price
\newcommand{\pindex}{p_{\text{index}}}       % FCI oracle index price
\newcommand{\Lact}{L_{\text{act}}}           % active liquidity
\newcommand{\Lnet}{\Delta L_{\text{net}}}    % net liquidity at tick crossing
\newcommand{\Ua}{U_a}                        % utilization ratio (borrowed/total)

% ── Borrow rate (JumpRate) ──────────────────────────────────────────────────
\newcommand{\rborrow}{r_{\text{borrow}}}     % borrow rate (PLP → underwriter)
\newcommand{\rfunding}{r_{\text{funding}}}   % funding rate (bidirectional)
\newcommand{\kink}{\kappa}                   % utilization kink point
\newcommand{\mlow}{m_1}                      % slope below kink
\newcommand{\mhigh}{m_2}                     % slope above kink (jump multiplier)

% ── Coordinate transform ────────────────────────────────────────────────────
% u = x²  (linearized risky reserve; x = virtual concentration exposure)
% y = numeraire reserve (USDC)
% ψ(u, y) = y - (p_l²/4)·u = L·C₁  (linear in (u,y), 1-homogeneous in L)


\title{Fee Concentration Insurance\\[4pt]
       \large Mathematical Specification}
\author{ThetaSwap}
\date{\today}

\begin{document}

\maketitle

\begin{abstract}
We specify an insurance mechanism against adverse fee concentration in
Uniswap V4 liquidity provision. The fee concentration index
$\AT = (\sum_k \thetapos_k \cdot x_k^2)^{1/2}$ measures realized fee
inequality across LP positions. The Ma--Crapis competitive null
$\ATnull = \sqrt{\ThetaSum/\PosCount^2}$ provides the equal-share baseline.
The concentration deviation $\DeltaPlus = \max(0,\, \AT - \ATnull)$
captures excess concentration above the competitive equilibrium and maps to
price space via the odds ratio $p = \DeltaPlus/(1-\DeltaPlus)$.

Econometric identification (quadratic deviation exit hazard model) establishes
an inverted-U relationship: for the pool we analyzed (ETH/USDC 30bps),
below the turning point $\DeltaStar \approx 0.09$, concentration is protective
(shelter effect); above it, concentration drives LP
exits\footnote{Capponi, Jia \& Zhu (2024) show that JIT liquidity
mechanically concentrates fees in short-lived positions, diluting passive LPs'
effective fee rate and pushing them toward their optimal exit threshold.}.
We find contingent empirical demand from passive liquidity providers for
insurance against adverse competition, and that the optimal payoff function
is quadratic. The turning point $\DeltaStar$ calibrates insurance demand.

The hook tracks three co-primary state variables---$\AT$ (fee concentration
index), $\ThetaSum$ (aggregate turnover rate $\sum_k \thetapos_k$), and
$\PosCount$ (active position count)---with $\DeltaPlus$ derived on the fly.
The mechanism is perpetual with jump-adjusted funding and USDC numeraire.
\end{abstract}

\tableofcontents

% payoff.tex — Fee concentration index, null hypothesis, deviation, price mapping
% Kontrol: prove_cfmm_feeconc_indexBoundedness, prove_cfmm_feeconc_nullConsistency

\section{Fee Concentration Index}\label{sec:fee-concentration-index}

\subsection{Primary Index: $\AT$}

The fee concentration index $\AT$ is a path-dependent, HHI-weighted measure of adverse competition
in liquidity provision. It is the \emph{primary} computed and tracked quantity:
\begin{equation}\label{eq:concentration-index}
  \AT = \Bigl(\sum_{k} \thetapos_k \cdot (x_k)^2 \Bigr)^{1/2}
\end{equation}
where the sum runs over all positions removed up to observation time $T$, and:
\begin{itemize}
  \item $x_k = \dfrac{\ell_k}{L(p^*)}$
        is the fee share captured by position $k$, where $\ell_k$ is the
        position's liquidity and $L(p^*)$ is the total liquidity in the
        active tick range;
  \item $\thetapos_k = \dfrac{1}{\Delta B_k}$
        is the sophistication weight, where $\Delta B_k = \text{block}(\text{removal}) - \text{block}(\text{creation})$
        is the position lifetime in blocks;
  \item $\thetapos_k = 1$ for JIT positions ($\Delta B_k = 1$ block: same-block add and remove).
\end{itemize}

The index satisfies $\AT \in [0,1]$.

\subsection{Co-Primary State Variables}

The hook tracks three state variables at each liquidity event:

\begin{definition}[Co-Primary State]\label{def:co-primary}
\begin{align}
  \AT &= \Bigl(\sum_{k} \thetapos_k \cdot x_k^2 \Bigr)^{1/2}
    && \text{fee concentration index} \label{eq:at-state}\\
  \ThetaSum &= \sum_{k} \thetapos_k
    && \text{theta sum (aggregate turnover rate)} \label{eq:theta-state}\\
  \PosCount &= |\{k : \text{position}_k \text{ active}\}|
    && \text{active position count} \label{eq:n-state}
\end{align}
All three are stored on-chain. The derived quantities $\ATnull$ and $\DeltaPlus$
are computed on the fly from these three values.
\end{definition}

The Ma--Crapis \cite{macrapis2024permissionless} symmetric equilibrium assumes all positions capture equal fee shares
$x_k = 1/\PosCount$. Under this null hypothesis:
\begin{equation}\label{eq:null-at}
  \ATnull = \sqrt{\frac{\ThetaSum}{\PosCount^2}}
\end{equation}

\begin{proposition}[Null as Lower Bound]\label{prop:null-lower-bound}
For any fee share distribution $\{x_k\}$ with $\sum_k x_k = 1$ and $x_k \geq 0$:
\begin{equation}
  \AT \geq \ATnull
\end{equation}
with equality if and only if $x_k = 1/\PosCount$ for all $k$.
\end{proposition}

\begin{proof}

 Under equal shares $x_k = 1/\PosCount$:
\[
  (\ATnull)^2 = \sum_k \thetapos_k \cdot \frac{1}{\PosCount^2}
             = \frac{\ThetaSum}{\PosCount^2}
\]

For arbitrary shares, by Jensen's inequality on the convex function $x \mapsto x^2$:
\[
  \sum_k \thetapos_k x_k^2 \geq \frac{\ThetaSum}{\PosCount^2}
  \quad \text{when } \sum_k x_k = 1
\]
since $\sum_k x_k^2 \geq 1/\PosCount$ by the constraint $\sum x_k = 1$
(this is the HHI lower bound). Multiplying each term by $\thetapos_k/\ThetaSum$
and using convexity preserves the inequality. Taking square roots gives $\AT \geq \ATnull$. \qedhere
\end{proof}

\subsection{Concentration Deviation: $\DeltaPlus$}\label{subsec:deviation}

The economically meaningful treatment variable is the excess concentration above the
competitive null:

\begin{definition}[Concentration Deviation]\label{def:delta-plus}
\begin{equation}\label{eq:delta-plus}
  \DeltaPlus = \max\bigl(0,\; \AT - \ATnull\bigr)
\end{equation}
\end{definition}

By \cref{prop:null-lower-bound}, $\DeltaPlus \geq 0$ always.
$\DeltaPlus = 0$ corresponds to the competitive equilibrium (all LPs earn equal fee shares).
$\DeltaPlus > 0$ indicates excess concentration: a few LPs capture disproportionate fee revenue.


\subsection{Price Mapping}\label{subsec:price-mapping}

\begin{definition}[Concentration Price]\label{def:price-mapping}
The price of fee concentration in numeraire (USDC) is the odds ratio:
\begin{equation}\label{eq:price-mapping}
  p = \frac{\DeltaPlus}{1 - \DeltaPlus}
\end{equation}
with $p \in [0, \infty)$ for $\DeltaPlus \in [0, 1)$.
The inverse mapping is $\DeltaPlus = p/(1+p)$.
\end{definition}

\begin{proposition}[Monotonicity]\label{prop:price-monotone}
$p(\DeltaPlus)$ is strictly increasing on $[0, 1)$.
\end{proposition}

\begin{proof}
$\dfrac{dp}{d\DeltaPlus} = \dfrac{1}{(1-\DeltaPlus)^2} > 0$ for all $\DeltaPlus < 1$. \qedhere
\end{proof}

\section{Econometric Calibration}\label{sec:econometric}

The econometric identification (see companion document \texttt{lp-insurance-demand.tex})
establishes the demand structure for fee concentration insurance.

\subsection{Quadratic Deviation Model}

The exit hazard model with quadratic treatment:
\begin{equation}\label{eq:quadratic-hazard}
  P(\text{exit}_{i,t} = 1) = \sigma\bigl(
    \beta_0 + \beta_1 \cdot \text{dev}_t + \beta_2 \cdot \text{dev}_t^2
    + \beta_3 \cdot \text{IL}_t + \beta_4 \cdot \ln(\text{age}_{i,t})
  \bigr)
\end{equation}
where $\text{dev}_t = \AT - \ATnull$ is the deviation from the competitive null.

\subsection{Key Results}

\begin{center}
\begin{tabular}{@{}lcccc@{}}
\toprule
Lag & $\beta_1$ (linear) & $p$ & $\beta_2$ (quadratic) & $p$ \\
\midrule
1 & $-23.18$ & $0.012^{**}$ & $+129.20$ & $0.030^{**}$ \\
2 & $-43.42$ & $0.016^{**}$ & $+226.92$ & $0.065^{*}$ \\
3 & $-32.44$ & $0.001^{***}$ & $+205.34$ & $0.004^{***}$ \\
\bottomrule
\end{tabular}
\end{center}

\subsection{Inverted-U and the Insurance Trigger}

The model produces an inverted-U with turning point:
\begin{equation}\label{eq:turning-point}
  \DeltaStar = -\frac{\beta_1}{2\beta_2} \approx 0.09
\end{equation}

\begin{itemize}
  \item \textbf{Below $\DeltaStar$ (shelter regime):} Mild fee concentration correlates with
    high volume. All LPs benefit from elevated fee revenue. $\partial P(\text{exit})/\partial \text{dev} < 0$.
  \item \textbf{Above $\DeltaStar$ (Capponi regime):} Extreme concentration means passive LPs'
    effective fee rate drops below their exit threshold.
    $\partial P(\text{exit})/\partial \text{dev} > 0$.
    Capponi \& Zhu (2024): $\partial p^*_U/\partial \varphi > 0$, so lower $\varphi \to$ exit.
\end{itemize}

The turning point $\DeltaStar \approx 0.09$ identifies when insurance demand activates.
This calibrates the protection mechanism regardless of architecture choice.

\section{Settled Design Properties}\label{sec:design-properties}

The following properties are confirmed by the ETH/USDC 30bps historical backtest
(41 days, 600 positions, 2025-12-05 to 2026-01-14) and hold independent of
implementation details.

\subsection{Backtest-Confirmed Properties}

\begin{enumerate}
  \item \textbf{Perpetual.} No fixed expiry. The instrument settles via a funding rate mechanism
    at liquidity events, following the framework of Angeris et al.\ \cite{angeris2022perpetuals}.

  \item \textbf{Jump-adjusted funding.} $\DeltaPlus$ can jump discontinuously when large positions
    mint or burn, changing $\PosCount$ and $\ThetaSum$ abruptly. The backtest confirms
    concentration jumps are rare but large (1 trigger day in 41 days, $\DeltaPlus$: $0 \to 0.158$).
    The funding rate accounts for this jump risk per the Primer on Perpetuals.

  \item \textbf{Discrete funding at liquidity events.} Funding is exchanged when
    \texttt{afterSwap}, \texttt{afterAddLiquidity}, or \texttt{afterRemoveLiquidity}
    fires in the underlying pool---matching the FCI update frequency.
    No keeper or per-block updates required.

  \item \textbf{USDC numeraire.} Insurance is denominated in USDC (the pool's fee token).

  \item \textbf{Oracle-assisted framework.} The FCI hook provides $\pindex = \DeltaPlus/(1-\DeltaPlus)$
    as the index price. Following Angeris et al.\ \cite{angeris2021replicating}, the oracle-assisted
    replication framework applies because the hook's own state
    ($\AT$, $\ThetaSum$, $\PosCount$) provides the index endogenously.

  \item \textbf{Monotone payoff.} The protection value is monotone non-decreasing in
    $p = \DeltaPlus/(1-\DeltaPlus)$. The inverted-U from econometrics calibrates
    \emph{demand} (where PLPs buy protection), not payoff shape.
\end{enumerate}

\subsection{Optimal Payoff Formula: Power-Squared Exit Payoff}\label{subsec:optimal-payoff}

The backtest compared seven candidate payoff formulas across the ETH/USDC 30bps
position dataset:

\begin{center}
\begin{tabular}{@{}llcc@{}}
\toprule
Mode & Formula & \% Better Off & Mean HV (\$) \\
\midrule
Trigger 1/N (INS-05) & $\text{payout} = \Res / N_{\text{insured}}$ & 0.0\% & $-107.67$ \\
Trigger premium-prop & $\text{payout} \propto \text{premium}_i$ & 2.7\% & $-95.16$ \\
Exit $\times\; \DeltaPlus$ & $\text{payout} = \premfactor \cdot f \cdot \DeltaPlus_{\max}$ & 4.2\% & $-88.03$ \\
Exit $\times\; p$ & $\text{payout} = \premfactor \cdot f \cdot p_{\max}$ & 8.1\% & $-71.22$ \\
\textbf{Exit $\times\; p^2$ ($\alpha=2$)} & $\textbf{payout} = \textbf{premium} \times ((p_{\max}/\pstar)^2 - 1)$ & \textbf{18.8\%} & $\mathbf{+23.14}$ \\
Exit $\times\; p^3$ ($\alpha=3$) & $\text{payout} = \text{premium} \times ((p_{\max}/\pstar)^3 - 1)$ & 12.1\% & $-14.52$ \\
\bottomrule
\end{tabular}
\end{center}
(All results at $R_0 = \$200\text{K}$ seed capital, $\premfactor = 3.30\%$, $\DeltaStar = 0.09$.)

\begin{definition}[Power-Squared Exit Payoff]\label{def:exit-payoff}
At position exit (\texttt{afterRemoveLiquidity}), PLP $i$ receives:
\begin{equation}\label{eq:exit-payoff}
  \text{payout}_i = \text{premium}_i \times \left[\left(\frac{\max_t\, p(t)}{\pstar}\right)^2 - 1\right]^{+}
\end{equation}
where $\text{premium}_i = \premfactor \times \text{lifetime\_fees}_i$,
$p(t) = \DeltaPlus(t)/(1-\DeltaPlus(t))$ is the concentration price at time $t$,
and $\max_t$ is taken over the position's lifetime.
\end{definition}

\subsection{Why Power-Squared Dominates}

\begin{enumerate}
  \item \textbf{No moral hazard.} The payout depends on the \emph{worst} (highest)
    concentration experienced during the position's lifetime. Staying in the pool
    during concentration spikes \emph{increases} the payout. Exiting early reduces it.
    This is the opposite of the premium-proportional moral hazard.

  \item \textbf{Tail sensitivity.} The $\alpha = 2$ exponent is the sweet spot for
    rare severe events at $\DeltaStar = 0.09$: sufficient convexity to reward
    extreme concentration without draining the reserve at moderate $\DeltaPlus$.
    Higher $\alpha$ (e.g., $\alpha = 3$) over-concentrates payouts on the single worst event;
    lower $\alpha$ (e.g., $\alpha = 1$) provides insufficient tail compensation.

  \item \textbf{Incentive alignment.} Backtest confirms positive correlation between
    position lifetime and hedge value: longer-lived positions benefit more than
    short-lived ones (the correct incentive structure for insurance).

  \item \textbf{Bootstrapping resolution.} Equal pro-rata (INS-05) yields 0\% of
    positions better off at any $\premfactor$ due to the bootstrapping problem:
    the mutual reserve is too small at trigger time. The exit payoff sidesteps
    bootstrapping entirely by settling at exit, not at trigger.
\end{enumerate}

\section{Architectural Design}\label{sec:architecture}

The architecture is a \textbf{Lendgine-style CFMM} following
Leifke \& Scott \cite{leifke2024powermaker}, adapted from ETH price to
fee concentration price $p = \DeltaPlus/(1-\DeltaPlus)$.

\subsection{Participants and Token Mapping}

\begin{definition}[Insurance CFMM Roles]\label{def:cfmm-roles}
\begin{itemize}
  \item \textbf{Underwriter (LP of insurance CFMM):} External actor who deposits USDC
    collateral into a tick range $[\pl, \pu]$. Earns borrow rate (premium stream).
    Commits to $-p^2$ payoff: loses collateral quadratically as concentration rises.

  \item \textbf{PLP (Borrower):} Uniswap V4 LP who borrows the insurance CFMM's LP
    share. Receives $+p^2$ protection payoff. Pays borrow rate from V4 accrued fees.

  \item \textbf{Token X (virtual):} Concentration exposure token---not a real ERC-20
    but an internal accounting unit representing units of concentration risk priced
    by the FCI oracle.

  \item \textbf{Token Y (real):} USDC---underwriter's collateral, PLP's premium currency.
\end{itemize}
\end{definition}

\subsection{LP Payoff (Underwriter)}

Within tick range $[\pl, \pu]$, the underwriter's portfolio value per unit liquidity:
\begin{equation}\label{eq:lp-payoff}
  V(p) = \begin{cases}
    \dfrac{\pu^2 - \pl^2}{\pl^2} & p \leq \pl \\[6pt]
    \dfrac{\pu^2 - p^2}{\pl^2}   & \pl \leq p \leq \pu \\[6pt]
    0                              & p \geq \pu
  \end{cases}
\end{equation}

\begin{proposition}[Concavity]\label{prop:lp-concavity}
$V(p)$ is concave: $V''(p) = -2/\pl^2 < 0$ in $[\pl, \pu]$.
\end{proposition}

\subsection{Borrower Payoff (PLP)}

The PLP borrows the LP share and receives the convex complement:
\begin{equation}\label{eq:plp-payoff}
  N(p) = V(\pl) - V(p) = \begin{cases}
    0                                                    & p \leq \pl \\[6pt]
    \left(\dfrac{p}{\pl}\right)^2 - 1                   & \pl \leq p \leq \pu \\[6pt]
    \left(\dfrac{\pu}{\pl}\right)^2 - 1                 & p \geq \pu
  \end{cases}
\end{equation}
This matches the backtest-confirmed payoff (\cref{def:exit-payoff}) with
$\pstar = \pl$ (the lowest active underwriter tick) and cap at $\pu$.

\subsection{Market-Implied Strike}

The strike $\pstar$ is \emph{not} a protocol parameter. It emerges from the market:
\begin{equation}\label{eq:market-strike}
  \pstar = \min\{\pl : \exists\; \text{underwriter with liquidity at } \pl\}
\end{equation}
Underwriters express their view on where concentration becomes dangerous by
choosing their tick range $[\pl, \pu]$. The aggregate liquidity profile across
ticks is the insurance pricing curve. As market conditions change, underwriters
reposition, and $\pstar$ adjusts naturally---no governance required.

\subsection{Connection to PowerMaker}

The architecture maps directly onto the PowerMaker \cite{leifke2024powermaker}
Lendgine construction:

\begin{center}
\begin{tabular}{@{}lll@{}}
\toprule
PowerMaker & ThetaSwap Insurance & Role \\
\midrule
ETH price $p$ & Concentration price $p = \DeltaPlus/(1-\DeltaPlus)$ & Underlying \\
LP (deposits into CFMM) & Underwriter (deposits USDC) & Sells convexity \\
Borrower (borrows LP share) & PLP (borrows LP share) & Buys convexity \\
$V(p) = 2pk - p^2$ & $V(p) = (\pu^2 - p^2)/\pl^2$ & LP payoff \\
$\varphi = R_1 - (k - \tfrac{1}{2}R_2)^2$ & $\psi = y - \tfrac{\pl^2}{4}u$ & Trading function \\
JumpRate borrow curve & JumpRate + $\premfactor$ base & Borrow rate \\
\bottomrule
\end{tabular}
\end{center}

% econometrics.tex — Structural Econometric Specification
% Fee Concentration Risk and LP Exit: Quadratic Deviation Model
% Source: lp-insurance-demand.pdf (2026-03-05)

\section{Methodology and Data}\label{sec:econometric-full}

This section presents the full econometric identification of insurance demand
from fee concentration risk.

The primary finding is that, for the ETH/USDC 30bps pool we study, there
exists a concentration threshold ($\DeltaStar \approx 0.09$) beyond which
JIT arrival crowds out passive liquidity providers. In pools where trading
volume is inelastic to liquidity depth \cite{capponi2024jit}, JIT positions
dilute passive LPs' effective fee rate without generating proportional
volume---shifting concentration from a signal of pool health to a driver of
LP exit. This regime shift is the empirical foundation of insurance demand.

The goal is to identify precisely where this transition occurs and to
translate the excess exit risk into a premium for fee concentration
derivatives. This motivates two questions:

\begin{enumerate}
  \item Does fee concentration risk ($\AT$) drive passive LP exits in
    Uniswap V3?
  \item At what concentration threshold does the exit effect activate, and
    can it be priced as an insurance premium?
\end{enumerate}

\subsubsection{Unit of Observation}

\textbf{Position-day panel.} One observation per position per day over a
41-day window (2025-12-05 to 2026-01-14). Binary outcome:
$\text{exit}_{i,t} \in \{0,1\}$. Treatment: fee concentration deviation from
the competitive null, measured at the pool-day level with lag.

\subsubsection{Outcome Variable}

Binary exit indicator:
\begin{equation}\label{eq:exit-indicator}
  Y_{i,t} = \begin{cases}
    1 & \text{if position } i \text{ burns on day } t \\
    0 & \text{otherwise}
  \end{cases}
\end{equation}

Panel size: 3{,}365 position-day observations, 597 exit events, 40 cluster
days, from 600 positions in the ETH/USDC 30bps pool.

%% ─── 2. The Fee Concentration Index ─────────────────────────────────────────

\subsection{Fee Concentration Index (Econometric Definition)}\label{subsec:econ-fci}

\begin{definition}[Fee Concentration Index]\label{def:econ-fci}
For a pool with $\PosCount$ active positions on day $t$, the fee concentration
index is:
\begin{equation}\label{eq:econ-at}
  \AT = \sqrt{\sum_{k=1}^{\PosCount} \thetapos_k \, x_k^2}
\end{equation}
where $x_k = \text{fee}_k / \sum_j \text{fee}_j$ is position $k$'s fee share
and $\thetapos_k = 1/\blocklife_k$ is its turnover rate.
\end{definition}

$\AT$ aggregates two dimensions of competitive pressure: \emph{who earns the
fees} ($x_k$) and \emph{how fast positions turn over} ($\thetapos_k$). High
$\AT$ means a few short-lived positions capture most fee revenue.

\subsubsection{The Equal-Share Null}

\begin{definition}[Ma--Crapis Null]\label{def:econ-null}
Under the Ma \& Crapis (2024) symmetric equilibrium with $\PosCount$ identical
LPs, each earns equal fee share $x_k = 1/\PosCount$:
\begin{equation}\label{eq:econ-null}
  \ATnull = \sqrt{\frac{\sum_k \thetapos_k}{\PosCount^2}}
\end{equation}
This is the competitive baseline---the value $\AT$ takes when fee
concentration is zero.
\end{definition}

\subsubsection{Concentration Deviation}

\begin{definition}[Concentration Deviation]\label{def:econ-deviation}
The treatment variable is the excess concentration above the null:
\begin{equation}\label{eq:econ-deviation}
  \Delta_t \equiv A_{T,t}^{\text{real}} - A_{T,t}^{1/N}
\end{equation}
computed from the \textbf{same} set of positions on each day.
$\Delta_t > 0$ indicates concentration above the competitive equilibrium;
$\Delta_t < 0$ indicates more equal fee distribution than the symmetric
benchmark.
\end{definition}

\subsubsection{Empirical Properties of $\Delta_t$}

Data source: Dune query 6783604, joining \texttt{Collect}, \texttt{Burn}, and
\texttt{Mint} events for the ETH/USDC 30bps pool
(\texttt{0x8ad599c3\ldots eb6e6d8}). Fee amounts isolated by subtracting Burn
principal from Collect amounts. Token ordering: token0 = USDC (6 dec),
token1 = WETH (18 dec).

\begin{center}
\begin{tabular}{@{}ll@{}}
\toprule
Statistic & Value \\
\midrule
Days with $\Delta > 0$ & 26/41 (63\%) \\
Mean $A_T^{\text{real}} / A_T^{1/N}$ & $2.65\times$ \\
Median deviation & $+0.0004$ \\
Range & $[-0.032,\; +0.158]$ \\
Positions per day & 13--65 \\
\bottomrule
\end{tabular}
\end{center}

The real $\AT$ exceeds the null on most days, confirming empirically that fee
shares are asymmetric---consistent with Aquilina et al.\ finding that 7\% of
LP addresses capture 80\% of fees.

%% ─── 3. Theoretical Foundations ─────────────────────────────────────────────

\subsection{Theoretical Foundations}\label{subsec:econ-theory}

Five papers independently validate the components of $\AT$ and the exit
mechanism.

\subsubsection{Capponi \& Zhu (2024): Optimal Exit Threshold}

\emph{``Optimal Exiting for Liquidity Provision in Constant Function Market Makers.''}

The LP solves an optimal stopping problem with exit thresholds $p_U^*$ and
$p_L^*$. The key comparative static:
\begin{equation}\label{eq:capponi-exit}
  \frac{\partial p_U^*}{\partial \phi} > 0
\end{equation}
A higher effective fee rate $\hat{\phi}$ raises the exit threshold (LP stays
longer). When $\AT$ rises, passive LP's effective fee rate
$\hat{\phi}_i = x_i \cdot \phi_{\text{pool}}$ drops because $x_i < 1/\PosCount$.
Lower $\hat{\phi} \Rightarrow$ lower $p_U^* \Rightarrow$ earlier exit.

\textbf{Implication for $\AT$:} $\AT$ enters the exit decision through
$\hat{\phi}$. This predicts $\beta_1 > 0$ for sufficiently high concentration.

\subsubsection{Capponi, Jia \& Zhu (2024): JIT Fee Dilution}

\emph{``The Paradox of Just-in-Time Liquidity in Decentralized Exchanges.''}

JIT liquidity creates fee share asymmetry: the passive LP's share is
$d_P/(d_P + d_J)$ where $d_J$ is JIT depth. This is exactly the $x_k$
component of $\AT$. JIT entry mechanically raises $\AT$ by concentrating fees
in short-lived positions.

\textbf{Implication for identification:} JIT activity is a natural instrument
for $\AT$. Our JIT control analysis (32 JIT positions, 5.3\% share) shows JIT
count has an independent positive effect on exits
($\beta_{\text{JIT}} = 0.177$, $p = 0.031$), but does not explain the $\AT$
effect---confirming that $\AT$ captures concentration \emph{beyond} JIT.

\subsubsection{Ma \& Crapis (2024): The Competitive Null}

\emph{``Decentralized Exchange Design: Permissionless Liquidity Provision.''}

$\PosCount$ symmetric LPs in free-entry equilibrium: each earns
$\Pi_i = (1/\PosCount^2) f B_D$. The equal-share outcome $x_k = 1/\PosCount$
is the \textbf{null hypothesis} of $\AT$. Our deviation
$\Delta_t = A_T^{\text{real}} - A_T^{1/N}$ tests departure from this
equilibrium directly.

\textbf{Implication for the model:} $\Delta_t = 0$ means the pool operates at
the competitive equilibrium. $\Delta_t > 0$ means fee concentration exceeds
what symmetric competition predicts. The deviation is the correct treatment
variable---not $\AT$ in levels.

\subsubsection{Aquilina, Foley, Gambacorta \& Krekel (2024): Empirical Validation}

\emph{``Decentralised Dealers? Examining Liquidity Provision in Decentralised
Exchanges.''} BIS Working Paper No.~1227.

Sophisticated LPs (7\% of addresses) provide 65--85\% of liquidity, use tick
ranges $3.4\times$ tighter, hold positions $7.5\times$ shorter, and capture
75--90\% of fees. This validates that $x_k$ in $\AT$ is highly skewed, not
symmetric.

Their profitability decomposition
$R_{\text{net}} = \text{FeeYield} + \text{IL} - \text{Gas}$ shows that passive
LPs earn negative returns when sophisticated LPs dominate---precisely when
$\Delta_t$ is large.

\subsubsection{Bichuch \& Feinstein (2024): Fee-as-Derivative Pricing}

\emph{``The Price of Liquidity: Implied Volatility of AMM Fees.''}

The LP value function is an optimal stopping problem with implied fee rate
$F(p_x, p_y) = p_y \cdot \ell(p_x/p_y)$ where $\ell$ is instantaneous LVR.
Theorem~4.2: a risk-neutral LP is indifferent across all stopping times. This
means \emph{any} deterministic exit driver must come from fee rate
heterogeneity across LPs---exactly what $\AT$ measures.

Their fixed-for-floating fee swap (Corollary~5.5) is structurally identical to
ThetaSwap's insurance product: the LP exchanges volatile realized fees for a
fixed premium.

\textbf{Implication:} $\AT$ measures the \emph{distribution} of the pool fee
rate across LPs. Bichuch \& Feinstein's framework prices the pool-level fee;
$\AT$ determines who receives it.

%% ─── 4. Economic Model ─────────────────────────────────────────────────────

\subsection{Economic Model}\label{subsec:econ-model}

\subsubsection{Environment and Actors}

Single Uniswap V3 pool (ETH/USDC 30bps). Two agent types:

\begin{enumerate}
  \item \textbf{Passive LP (PLP):} Commits liquidity for extended periods.
    Observes pool state (reserves, price, accumulated fees). Does \emph{not}
    observe $\AT$ directly but reacts to its symptoms (lower fee revenue).
    Decision: when to exit.

  \item \textbf{JIT/Sophisticated LP:} Not explicitly modeled as strategic.
    Enters as a latent force that raises $\AT$ by capturing disproportionate
    fee shares. Characterized by short blocklife and tight tick ranges
    (Aquilina et al.).
\end{enumerate}

\subsubsection{Equilibrium Concept}

Competitive (price-taking) per Capponi \& Zhu (2024). The PLP takes fee rate,
price process, and concentration risk as given and solves an optimal stopping
problem.

\textbf{Assumption} (Price-Taking LP). \emph{The PLP solves:}
\begin{equation}\label{eq:plp-stopping}
  \max_{\tau}\; \mathbb{E}\!\left[
    \int_0^{\tau} e^{-rt} \hat{\phi}_i \, dF_t - dL_t
    \;\middle|\; \mathcal{F}_0
  \right]
\end{equation}
where $\hat{\phi}_i = x_i \cdot \phi_{\text{pool}}$ is the LP's effective fee
rate given fee share $x_i$, $F_t$ is cumulative fee revenue, and $L_t$ is
impermanent loss.

%% ─── 5. Stochastic Model ───────────────────────────────────────────────────

\subsection{Stochastic Model}\label{subsec:econ-stochastic}

\subsubsection{Unobserved Heterogeneity ($\eta$)}

\begin{enumerate}
  \item \textbf{LP risk aversion} ($\eta_i^{\text{risk}}$): heterogeneous
    preferences. Some PLPs tolerate more concentration. NFT Position Manager
    masks identity.

  \item \textbf{Outside options} ($\eta_i^{\text{outside}}$): competing pool
    returns, staking yields, CEX opportunities.
\end{enumerate}

Bias direction: \textbf{toward zero} (conservative). Risk-averse LPs avoid
entering during high-$\AT$ periods \emph{and} exit faster, attenuating
$\beta$.

\subsubsection{Agent Uncertainty}

Future price process only ($u_t \sim$ GBM, per Capponi \& Zhu). Drives IL
realization.

\subsubsection{Implied Error Structure}

\begin{equation}\label{eq:error-structure}
  \varepsilon_i =
    \underbrace{\eta_i^{\text{risk}} + \eta_i^{\text{outside}}}_{\text{unobserved heterogeneity}}
    + \underbrace{\nu_i}_{\text{measurement error}}
\end{equation}

%% ─── 6. Estimation Strategy ────────────────────────────────────────────────

\subsection{Estimation Strategy}\label{subsec:econ-estimation}

\subsubsection{Model Progression}

We estimate three nested specifications to motivate the quadratic model.

\paragraph{Specification 1: Linear (Point Lag).}
\begin{equation}\label{eq:spec-linear-point}
  P(\text{exit}_{i,t} = 1) = \sigma\bigl(
    \beta_0 + \beta_1 A_{T,t-\ell} + \beta_2 \,\text{IL}_t
    + \beta_3 \log(\text{age}_{i,t})
  \bigr)
\end{equation}

\textbf{Result 1} (Linear A\_T). $\beta_1 = -3.93$, cluster $p = 0.072$.
Marginally significant but \textbf{wrong sign}---higher concentration
associated with fewer exits. Robust across lags 1--7 (all negative, 4/5
significant).

\paragraph{Specification 2: Linear Deviation from $1/\PosCount$.}
\begin{equation}\label{eq:spec-linear-dev}
  P(\text{exit}_{i,t} = 1) = \sigma\bigl(
    \beta_0 + \beta_1 \Delta_{t-\ell} + \beta_2 \,\text{IL}_t
    + \beta_3 \log(\text{age}_{i,t})
  \bigr)
\end{equation}
where $\Delta_t = A_T^{\text{real}} - A_T^{1/N}$ from the same position set.

\textbf{Result 2} (Linear Deviation). $\beta_1 = -5.44$, cluster $p = 0.109$.
Still negative. But \textbf{dose-response analysis reveals non-monotonicity}:

\begin{center}
\begin{tabular}{@{}llll@{}}
\toprule
Quartile & Mean $\Delta$ & Exit Rate & Pattern \\
\midrule
Q1 (below null) & $-0.009$ & 17.75\% & baseline \\
Q2 (at null) & $-0.000$ & 17.87\% & baseline \\
\textbf{Q3 (mild conc.)} & $\mathbf{+0.002}$ & \textbf{20.56\%} & \textbf{exits increase} \\
Q4 (extreme conc.) & $+0.025$ & 14.37\% & exits decrease \\
\bottomrule
\end{tabular}
\end{center}

Q3 shows the predicted Capponi \& Zhu pattern. Q4 reverses---motivating the
quadratic specification.

\paragraph{Specification 3: Quadratic Deviation (Primary Model).}
\begin{equation}\label{eq:quadratic-hazard-full}
  \boxed{P(\text{exit}_{i,t} = 1) = \sigma\bigl(
    \beta_0 + \beta_1 \Delta_{t-\ell} + \beta_2 \Delta_{t-\ell}^2
    + \beta_3 \,\text{IL}_t + \beta_4 \log(\text{age}_{i,t})
  \bigr)}
\end{equation}

This is the \textbf{primary specification}. The quadratic term captures the
inverted-U: shelter at low concentration, Capponi exit at high concentration.

\subsubsection{Estimation Method}

Logit MLE via Newton--Raphson (IRLS) with \textbf{day-clustered sandwich
standard errors} (41 clusters):
\begin{equation}\label{eq:clustered-var}
  \widehat{\text{Var}}(\hat{\beta}) = H^{-1}
    \left(\sum_{g=1}^{G} S_g S_g'\right)
    H^{-1} \cdot \frac{G}{G-1} \cdot \frac{n}{n-k}
\end{equation}
where $S_g = \sum_{i \in g} (y_i - \hat{p}_i) x_i$ is the cluster-level score
and $H$ is the Hessian.

Implemented in JAX 0.9.1 CPU. Code: \texttt{econometrics/hazard.py}.

%% ─── 7. Results ─────────────────────────────────────────────────────────────

\subsection{Results}\label{subsec:econ-results}

\subsubsection{Main Result: Quadratic Deviation Model}

\begin{center}
\begin{tabular}{@{}cccccl@{}}
\toprule
Lag & $\hat{\beta}_1$ (linear) & $p$ & $\hat{\beta}_2$ (quadratic) & $p$ & Turning point \\
\midrule
\textbf{1} & $\mathbf{-23.18}$ & $\mathbf{0.012}$ & $\mathbf{+129.20}$ & $\mathbf{0.030}$ & $\mathbf{0.090}$ \\
\textit{2} & \textit{$-43.42$} & \textit{0.016} & \textit{$+226.92$} & \textit{0.065} & \textit{0.096} \\
\textbf{3} & $\mathbf{-32.44}$ & $\mathbf{0.001}$ & $\mathbf{+205.34}$ & $\mathbf{0.004}$ & $\mathbf{0.079}$ \\
5 & $+3.81$ & 0.830 & $-76.31$ & 0.470 & --- \\
7 & $-5.55$ & 0.730 & $-16.46$ & 0.900 & --- \\
\bottomrule
\end{tabular}
\end{center}

Both terms of the quadratic deviation model are statistically significant at
the economically relevant lag horizon (1--3 days). The turning point clusters
at $\DeltaStar \approx 0.08$--$0.10$. The signal dissipates at lags 5--7,
consistent with PLPs responding to recent concentration shocks, not stale ones.

\subsubsection{Full Lag-1 Coefficients}

\begin{center}
\begin{tabular}{@{}lcrl@{}}
\toprule
Parameter & Estimate & Cluster SE & Interpretation \\
\midrule
$\beta_0$ (intercept) & varies & --- & baseline exit probability \\
$\beta_1$ ($\Delta$) & $-23.18$ & 9.21 & shelter below turning point \\
$\beta_2$ ($\Delta^2$) & $+129.20$ & 59.69 & exit above turning point \\
$\beta_3$ (IL) & $+13.52$ & --- & higher IL $\to$ more exits \\
$\beta_4$ (log age) & $-0.42$ & --- & older positions less likely to exit \\
\bottomrule
\end{tabular}
\end{center}

$n = 3{,}365$; exits $= 597$; clusters $= 40$; pseudo-$R^2 = 0.350$;
mean $P(\text{exit}) = 0.177$.

%% ─── 8. Economic Interpretation ─────────────────────────────────────────────

\subsection{Economic Interpretation}\label{subsec:econ-interpretation}

\subsubsection{The Inverted-U: Two Regimes}

The quadratic model identifies two regimes separated by the turning point
$\DeltaStar \approx 0.09$:

\begin{enumerate}
  \item \textbf{Shelter regime} ($\Delta < \DeltaStar$): Mild fee concentration
    correlates with high trading volume. All LPs benefit from elevated fee
    revenue, even if shares are unequal. The marginal effect
    $\partial P / \partial \Delta < 0$: more concentration, fewer exits.

  \item \textbf{Capponi regime} ($\Delta > \DeltaStar$): Extreme concentration
    means passive LPs' effective fee rate $\hat{\phi}_i$ drops below their exit
    threshold. Per Capponi \& Zhu:
    $\partial p_U^* / \partial \phi > 0$, so lower $\phi \Rightarrow$ lower
    $p_U^* \Rightarrow$ earlier exit. The marginal effect flips:
    $\partial P / \partial \Delta > 0$.
\end{enumerate}

\subsubsection{Marginal Effect at the Turning Point}

The marginal effect of the quadratic model is:
\begin{equation}\label{eq:marginal-effect}
  \frac{\partial P(\text{exit})}{\partial \Delta}
    = (\beta_1 + 2\beta_2 \Delta) \cdot \bar{P}(1 - \bar{P})
\end{equation}

At the turning point $\DeltaStar = -\beta_1/(2\beta_2) = 0.090$:
\[
  \left.\frac{\partial P}{\partial \Delta}\right|_{\DeltaStar} = 0
\]

For $\Delta = 0.15$ (well into the Capponi regime):
\begin{align}
  \beta_1 + 2\beta_2(0.15) &= -23.18 + 2(129.20)(0.15) = 15.58
    \label{eq:marginal-015} \\
  \frac{\partial P}{\partial \Delta} &= 15.58 \times 0.177 \times 0.823
    = 2.27 \label{eq:marginal-value}
\end{align}

A 0.01 increase in $\Delta$ at this level raises exit probability by
${\sim}2.3$ percentage points.

\subsubsection{Insurance Pricing Translation}

For a passive LP experiencing $\Delta = 0.15$ (concentration 15\% above the
competitive null):

\begin{center}
\begin{tabular}{@{}ll@{}}
\toprule
Quantity & Value \\
\midrule
$\Delta P(\text{exit})$ per 0.01 $\Delta$ & $+2.3$ pp \\
Average remaining position hours & 48 \\
Hours of fee revenue at risk & $0.023 \times 48 = 1.1$ hrs \\
Fee revenue per hour (pool avg) & \$100 \\
Implied premium per position & \$110 \\
\bottomrule
\end{tabular}
\end{center}

This is the maximum willingness-to-pay for ThetaSwap's insurance: a PLP would
pay up to \$110 to eliminate the excess exit risk from concentration above the
turning point.

%% ─── 9. Specification Tests ─────────────────────────────────────────────────

\subsection{Specification Tests}\label{subsec:econ-tests}

\subsubsection{Test 1: Sign Restriction (Capponi \& Zhu Prediction)}

\textbf{Hypothesis 1.} $\beta_2 > 0$: the quadratic term is positive,
indicating that extreme concentration drives exits.

\textbf{Result: PASS.} $\hat{\beta}_2 = +129.20$, $p = 0.030$ at lag~1;
$+205.34$, $p = 0.004$ at lag~3. The Capponi exit mechanism is confirmed in
the high-concentration regime.

\subsubsection{Test 2: Dose-Response Monotonicity}

\textbf{Hypothesis 2.} Exit rates increase with $\Delta$ in the upper
quartiles.

\textbf{Result: PASS (inverted-U).} Q1--Q2 stable (17.8\%), Q3 peak (20.6\%),
Q4 drops (14.4\%). The non-monotonicity motivates the quadratic and is
captured by the model.

\subsubsection{Test 3: Lag Robustness}

\textbf{Hypothesis 3.} Both $\beta_1$ and $\beta_2$ are significant for at
least 3 of 5 lag windows.

\textbf{Result: PASS.} Lags 1, 2, 3 all have significant quadratic terms
($p < 0.07$). The turning point is stable at 0.08--0.10. Signal dissipates at
lags 5--7, consistent with short-memory PLP reaction.

\subsubsection{Test 4: IL Orthogonality}

\textbf{Hypothesis 4.} $\beta_1, \beta_2$ persist with and without the IL
control.

\textbf{Result: PASS.} IL coefficient ($\beta_3 = 13.52$) is positive
(higher IL $\to$ more exits) but does not absorb the concentration effect.
The quadratic structure is present in the $\AT$-only model as well.

\subsubsection{Test 5: JIT Confounding Control}

\textbf{Hypothesis 5.} The concentration effect is not an artifact of JIT
activity.

\textbf{Result: PASS.} Adding a JIT proxy (positions with blocklife $< 100$
blocks) as a control:
\begin{itemize}
  \item JIT count has independent positive effect on exits
    ($\beta_{\text{JIT}} = 0.177$, $p = 0.031$)
  \item The $\AT$ effect becomes \emph{more} significant
    ($p: 0.072 \to 0.022$) with JIT controlled
  \item Confirms: $\AT$ captures concentration \emph{beyond} JIT
\end{itemize}

\subsubsection{Test 6: Alternative Treatment Timing}

Three treatment constructions tested:

\begin{center}
\begin{tabular}{@{}llll@{}}
\toprule
Treatment & $\beta_1$ (linear) & $p$ & Finding \\
\midrule
Point A\_T (real, lag 1) & $-3.93$ & 0.072 & Negative (shelter) \\
Lifetime mean A\_T & $-7.60$ & 0.444 & Negative, insignificant \\
Deviation from $1/\PosCount$ & $-5.44$ & 0.109 & Negative (Q3 reveals non-linearity) \\
\textbf{Quadratic deviation} & \multicolumn{2}{l}{\textbf{See Result 3}} & \textbf{Inverted-U confirmed} \\
\bottomrule
\end{tabular}
\end{center}

%% ─── 10. Robustness: Five Linear Specifications ─────────────────────────────

\subsection{Robustness: Five Linear Specifications}\label{subsec:econ-robustness}

Before arriving at the quadratic model, we tested five linear specifications
that progressively refined the treatment variable. All failed to produce a
positive $\beta_1$, but each contributed to understanding the data-generating
process.

\subsubsection{Specification A: Point Real A\_T}

Treatment $= A_{T,t-\ell}^{\text{real}}$ (single day's concentration).

\begin{center}
\begin{tabular}{@{}ccccc@{}}
\toprule
Lag & $\hat{\beta}_1$ & Cluster SE & $p$ & Sig \\
\midrule
1 & $-3.928$ & 2.186 & 0.072 & * \\
2 & $-8.783$ & 2.227 & 0.0001 & *** \\
3 & $-0.414$ & 1.131 & 0.714 & \\
5 & $-6.523$ & 1.650 & 0.0001 & *** \\
7 & $-5.625$ & 2.261 & 0.013 & ** \\
\bottomrule
\end{tabular}
\end{center}

\textbf{Finding:} Consistently negative. A naive interpretation is that
concentration \emph{protects} LPs. But the proxy A\_T ($x_k = 1/\PosCount$)
showed the same pattern ($\beta_1 = -4.90$, $p = 0.15$), suggesting the level
of $\AT$ confounds pool health (volume) with concentration pressure.

\subsubsection{Specification B: Lifetime Mean A\_T}

Treatment $=$ mean $A_T^{\text{real}}$ over $[\text{mint}_i, t - \ell]$.
Motivated by the hypothesis that passive LPs react to accumulated exposure,
not point shocks.

\begin{center}
\begin{tabular}{@{}ccccc@{}}
\toprule
Lag & $\hat{\beta}_1$ & Cluster SE & $p$ & Sig \\
\midrule
1 & $-7.595$ & 9.913 & 0.444 & \\
3 & $-11.603$ & 7.959 & 0.145 & \\
5 & $-14.977$ & 8.021 & 0.062 & * \\
7 & $-3.042$ & 6.829 & 0.656 & \\
\bottomrule
\end{tabular}
\end{center}

\textbf{Finding:} Amplifies the negative effect but with much larger SEs
(reduced variation in running means). Marginal significance only at lag~5.

\subsubsection{Specification C: Deviation from $1/\PosCount$ (Wrong Null)}

First attempt at the deviation treatment, using the old proxy A\_T from a
\emph{different} position set (Q4 burns) as the null. The real and null A\_T
were computed from different populations---an apples-to-oranges comparison.

\textbf{Finding:} Real A\_T was $10$--$100\times$ smaller than the proxy null
(different position sets, different blocklives). Deviation was negative for
40/41 days. Only lag-1 showed positive $\beta_1$ ($+1.92$, $p = 0.43$).
\textbf{Discarded} after identifying the data mismatch.

\subsubsection{Specification D: Corrected Deviation (Same Position Set)}

Null A\_T recomputed from the \emph{same} Dune query (same positions, same
blocklives). Now 26/41 days have $\Delta > 0$, mean ratio $= 2.65\times$.

\textbf{Finding:} Still negative $\beta_1$ in linear model, but dose-response
revealed the Q3 peak (20.56\%) and Q4 dip (14.37\%) that motivated the
quadratic specification.

\subsubsection{Specification E: JIT-Controlled Linear}

Added daily JIT count (positions with blocklife $< 100$ blocks) as control.

\textbf{Finding:} $\beta_1$ became \emph{more} negative and more significant
($p: 0.072 \to 0.022$). JIT itself predicts exits but does not confound the
$\AT$ effect.

%% ─── 11. Literature Alignment ───────────────────────────────────────────────

\subsection{Literature Alignment}\label{subsec:econ-literature}

Each component of the quadratic deviation model is validated by at least one
paper:

\begin{center}
\begin{tabular}{@{}p{3.5cm}p{4cm}p{5.5cm}@{}}
\toprule
Paper & Component Validated & Connection to Our Result \\
\midrule
Capponi \& Zhu (2024) &
  Exit threshold: $\partial p_U^*/\partial\phi > 0$ &
  Confirmed in the Capponi regime ($\Delta > 0.09$):
  $\beta_2 > 0$ means extreme concentration drives exits \\[6pt]
Capponi, Jia \& Zhu (2024) &
  $x_k$ captures JIT fee dilution &
  JIT control shows $\AT$ effect is independent of JIT confounding;
  $x_k$ reflects broader concentration than just JIT \\[6pt]
Ma \& Crapis (2024) &
  $x_k = 1/\PosCount$ is the competitive null &
  Deviation $\Delta = A_T^{\text{real}} - A_T^{1/N}$ is the correct treatment;
  levels confound pool health \\[6pt]
Aquilina et al.\ (2024) &
  7\% of LPs capture 80\% of fees &
  Real $\AT$ is $2.65\times$ the null on average, confirming massive fee share
  asymmetry \\[6pt]
Bichuch \& Feinstein (2024) &
  Risk-neutral LP indifferent on stopping times &
  Any exit driver must come from fee \emph{heterogeneity} across LPs---exactly
  what $\Delta$ measures \\
\bottomrule
\end{tabular}
\end{center}

%% ─── 12. Connection to ThetaSwap ────────────────────────────────────────────

\subsection{Connection to ThetaSwap}\label{subsec:econ-thetaswap}

\begin{enumerate}
  \item \textbf{Insurance trigger:} The turning point $\DeltaStar \approx 0.09$
    defines when insurance becomes valuable. When concentration exceeds 9\%
    above the Ma--Crapis equilibrium, passive LP exit risk increases.
    ThetaSwap monitors $\Delta_t$ in real time and activates payouts when
    $\Delta > \DeltaStar$.

  \item \textbf{Premium calibration:} The quadratic marginal effect at any
    $\Delta$ translates directly to an insurance premium via
    \cref{eq:marginal-effect}. At $\Delta = 0.15$: \$110 per position per
    event.

  \item \textbf{Distinct from IL hedging:} The IL coefficient ($\beta_3$) is
    positive but does not absorb the concentration effect. $\AT$ captures a
    risk orthogonal to price-movement risk, confirming that ThetaSwap
    addresses a gap unmet by options or IL protection products.

  \item \textbf{Product structure:} Per Bichuch \& Feinstein's
    fixed-for-floating fee swap framework, ThetaSwap is a derivative that pays
    the PLP when realized fee share $x_i$ falls below $1/\PosCount$ by more
    than $\DeltaStar$. The premium is the implied volatility of fee shares.

  \item \textbf{Market sizing:} 63\% of days in our sample have
    $\Delta > 0$. Of these, days exceeding $\DeltaStar$ represent the
    insurance-triggering events. Aggregating the implied premium across 600
    positions gives the minimum viable insurance pool size.
\end{enumerate}

%% ─── 13. Data and Implementation ────────────────────────────────────────────

\subsection{Data and Implementation}\label{subsec:econ-data}

\begin{center}
\begin{tabular}{@{}llll@{}}
\toprule
Query & Description & Rows & Cost \\
\midrule
Q4v2 & Position-level burns, blocklife & 600 & 0.18 cr \\
Q5 & Daily IL proxy & 41 & 0.14 cr \\
Q6 & Per-position Collect fees + blocklife (real $\AT$ and null $\AT$ per day)
  & 41 (aggregated) & 3.3 cr \\
\bottomrule
\end{tabular}
\end{center}

Total cost: ${\sim}3.6$ credits.

\textbf{Pool:} ETH/USDC 30bps (\texttt{0x8ad599c3a0ff1de082011efddc58f1908eb6e6d8}).\\
\textbf{Window:} 2025-12-05 to 2026-01-14 (41 days).\\
\textbf{Implementation:} Python 3.11, JAX 0.9.1 CPU, \texttt{@functional-python}
patterns (frozen dataclasses, pure functions).
Code in \texttt{econometrics/\{types,data,ingest,hazard\}.py}. 52 unit tests
passing.

Key functions:
\begin{itemize}
  \item \texttt{build\_exit\_panel\_deviation()}: constructs position-day panel
    with $\Delta$ as treatment
  \item \texttt{logit\_mle\_quadratic()}: Newton--Raphson logit with
    $\Delta + \Delta^2$ and clustered SEs
  \item \texttt{marginal\_effect\_at\_means()}: translates $\beta$ to insurance
    premium
\end{itemize}

%% ─── 14. References ─────────────────────────────────────────────────────────

\subsection{Econometric References}

\begin{itemize}
  \item Aquilina, M., Foley, S., Gambacorta, L., \& Krekel, M.\ (2024).
    ``Decentralised Dealers? Examining Liquidity Provision in Decentralised
    Exchanges.'' \emph{BIS Working Paper} No.~1227.
  \item Bichuch, M.\ \& Feinstein, Z.\ (2024).
    ``The Price of Liquidity: Implied Volatility of AMM Fees.''
  \item Capponi, A.\ \& Zhu, B.\ (2024).
    ``Optimal Exiting for Liquidity Provision in Constant Function Market Makers.''
  \item Capponi, A., Jia, R., \& Zhu, B.\ (2024).
    ``The Paradox of Just-in-Time Liquidity in Decentralized Exchanges:
    More Providers Can Lead to Less Liquidity.''
  \item Ma, J.\ \& Crapis, H.\ (2024).
    ``Decentralized Exchange Design: Permissionless Liquidity Provision.''
  \item Reiss, P.C.\ \& Wolak, F.A.\ (2007).
    ``Structural Econometric Modeling: Rationales and Examples from
    Industrial Organization.'' \emph{Handbook of Econometrics}, Vol.~6A, Ch.~64.
\end{itemize}

% trading-function.tex — Linearized power² trading function derivation
% Kontrol: prove_cfmm_tradingFunction*, prove_cfmm_coordinateTransform*

\section{Trading Function}\label{sec:trading-function}

\subsection{Reserve Derivation}

Starting from the LP payoff $V(p) = (\pu^2 - p^2)/\pl^2$ for $p \in [\pl, \pu]$,
we derive reserves via the Angeris replication theorem \cite{angeris2021replicating}.

\begin{definition}[Reserves]\label{def:reserves}
The risky (virtual concentration) and numeraire (USDC) reserves per unit liquidity:
\begin{align}
  x(p) &= -V'(p) = \frac{2p}{\pl^2} \label{eq:risky-reserve} \\
  y(p) &= V(p) + p \cdot x(p) = \frac{\pu^2 + p^2}{\pl^2} \label{eq:numeraire-reserve}
\end{align}
for $p \in [\pl, \pu]$, with boundary extensions:
$x = 2\pl/\pl^2$ and $y = (\pu^2 + \pl^2)/\pl^2$ for $p \leq \pl$;
$x = 2\pu/\pl^2$ and $y = 2\pu^2/\pl^2$ for $p \geq \pu$.
\end{definition}

\begin{proposition}[Reserve Verification]\label{prop:reserve-verify}
$V(p) = y(p) - p \cdot x(p)$ holds exactly:
\[
  \frac{\pu^2 + p^2}{\pl^2} - p \cdot \frac{2p}{\pl^2}
  = \frac{\pu^2 + p^2 - 2p^2}{\pl^2}
  = \frac{\pu^2 - p^2}{\pl^2}
  = V(p) \qquad\checkmark
\]
\end{proposition}

\subsection{Coordinate Transform: $u = x^2$}\label{subsec:coordinate-transform}

The natural trading function $y = \pu^2/\pl^2 + (\pl^2/4)x^2$ is quadratic in $x$.
This creates a 2-homogeneous invariant, incompatible with O(1) tick-based
liquidity aggregation (Uniswap V3-style $\Lact \mathrel{+}= \Lnet$ at tick crossings).

\begin{definition}[Linearization]\label{def:linearization}
Define the transformed risky reserve $u = x^2$. Then:
\begin{equation}\label{eq:linearized-invariant}
  \boxed{\psi(u, y) = y - \frac{\pl^2}{4}\, u}
\end{equation}
with invariant constant $C_1 = \pu^2/\pl^2$ per unit liquidity.
\end{definition}

\begin{theorem}[Linearity and 1-Homogeneity]\label{thm:linearity}
$\psi(u, y)$ is linear in $(u, y)$. With liquidity $L$:
\[
  u_{\text{total}} = L \cdot \tilde{u}(p), \quad
  y_{\text{total}} = L \cdot \tilde{y}(p)
\]
where $\tilde{u}(p) = 4p^2/\pl^4$ and $\tilde{y}(p) = (\pu^2 + p^2)/\pl^2$
are per-unit-liquidity reserves. The invariant scales as:
\[
  \psi(u_{\text{total}}, y_{\text{total}}) = L \cdot C_1
\]
which is 1-homogeneous in $L$.
\end{theorem}

\begin{proof}
$\psi(L\tilde{u}, L\tilde{y}) = L\tilde{y} - (\pl^2/4) \cdot L\tilde{u}
= L(\tilde{y} - (\pl^2/4)\tilde{u}) = L \cdot C_1$. \qedhere
\end{proof}

\begin{remark}[Concavity]
The Hessian of $\psi$ has eigenvalues $-\pl^2/2$ and $0$, both $\leq 0$.
Hence $\psi$ is concave, and the trading set
$S = \{(u, y) : \psi(u, y) \geq L C_1\}$ is convex
(Geometry paper \cite{angeris2023geometry}, Theorem~1).
\end{remark}

\subsection{Spot Price}

\begin{definition}[Spot Price from Reserves]\label{def:spot-price}
\begin{equation}\label{eq:spot-price}
  p = \frac{\pl^2}{2}\sqrt{\frac{u}{L}}
  = \frac{\pl^2}{2} \cdot \frac{x}{L}
  \quad\text{(in } x \text{ space: } p = \pl^2 x / (2L)\text{)}
\end{equation}
where the second equality uses $x = \sqrt{u}$.
\end{definition}

\begin{proof}[Verification]
From the trading function, $-\partial\psi/\partial u \big/ \partial\psi/\partial y
= -(-\pl^2/4) / 1 = \pl^2/4$. This is the marginal price in $(u, y)$ space.
In $(x, y)$ space, $du = 2x\, dx$, so the marginal price is
$dy/dx = (\pl^2/4) \cdot 2x = \pl^2 x/2 = p$, matching the reserve
derivation $x(p) = 2p/\pl^2 \Leftrightarrow p = \pl^2 x/2$. \qedhere
\end{proof}

\subsection{Swap Mechanics}

A swap in the insurance CFMM operates in $(u, y)$ space internally, with
$x \leftrightarrow u$ conversion at the boundary.

\begin{definition}[Swap]\label{def:swap}
Trader sells $\Delta x$ units of concentration exposure:
\begin{align}
  \Delta u &= (x + \Delta x)^2 - x^2 = 2x\,\Delta x + (\Delta x)^2 \label{eq:delta-u} \\
  \Delta y &= \frac{\pl^2}{4}\, \Delta u \label{eq:delta-y}
\end{align}
The effective price $\Delta y / \Delta x = (\pl^2/4)(2x + \Delta x)$ depends on
the current state $x$ and trade size $\Delta x$, providing natural slippage
despite the internal invariant being linear.
\end{definition}

\begin{remark}[O(1) Complexity]
Within a tick: the swap is a single multiplication (\cref{eq:delta-y}) plus
one square root for $x = \sqrt{u}$ conversion (\texttt{FixedPointMathLib.sqrt}).
Tick crossing: $\Lact \mathrel{+}= \Lnet$, a single addition.
Total per-swap cost: O(ticks crossed), same as Uniswap V3.
\end{remark}

\subsection{Liquidity Provisioning}

\begin{definition}[Single-Sided USDC Deposit]\label{def:single-sided}
An underwriter deposits $\Delta y$ USDC at current price $p$:
\begin{align}
  \Delta L &= \frac{\Delta y \cdot \pl^2}{\pu^2 + p^2} \label{eq:mint-L} \\
  \Delta u &= \Delta L \cdot \frac{4p^2}{\pl^4}
             = \frac{4p^2 \cdot \Delta y}{\pl^2(\pu^2 + p^2)} \label{eq:mint-u}
\end{align}
The virtual $\Delta u$ is computed, not deposited. Only USDC enters the contract.
\end{definition}

\begin{proposition}[Invariant Preservation on Mint]\label{prop:mint-invariant}
After minting: $\psi(u + \Delta u, y + \Delta y) = (L + \Delta L) \cdot C_1$.
\end{proposition}

\begin{proof}
$\psi(u + \Delta u, y + \Delta y) = (y + \Delta y) - (\pl^2/4)(u + \Delta u)
= \psi(u, y) + \Delta y - (\pl^2/4)\Delta u
= L C_1 + \Delta L \cdot \tilde{y} - (\pl^2/4)\Delta L \cdot \tilde{u}
= L C_1 + \Delta L \cdot C_1 = (L + \Delta L) C_1$. \qedhere
\end{proof}

% reserves.tex — Reserve curves and coordinate transform verification
% Kontrol: prove_cfmm_reserve*

\section{Reserve Analysis}\label{sec:reserves}

This section verifies the reserve curves derived in \cref{sec:trading-function}
and establishes boundary behavior and numerical properties.

\subsection{Reserve Curves in Price Space}

Per unit liquidity ($L = 1$), with $p \in [\pl, \pu]$:

\begin{center}
\begin{tabular}{@{}lccc@{}}
\toprule
& $p = \pl$ & $p \in (\pl, \pu)$ & $p = \pu$ \\
\midrule
$x(p) = 2p/\pl^2$ & $2/\pl$ & $2p/\pl^2$ & $2\pu/\pl^2$ \\
$u(p) = x^2 = 4p^2/\pl^4$ & $4/\pl^2$ & $4p^2/\pl^4$ & $4\pu^2/\pl^4$ \\
$y(p) = (\pu^2+p^2)/\pl^2$ & $(\pu^2+\pl^2)/\pl^2$ & $(\pu^2+p^2)/\pl^2$ & $2\pu^2/\pl^2$ \\
$V(p) = (\pu^2-p^2)/\pl^2$ & $(\pu^2-\pl^2)/\pl^2$ & $(\pu^2-p^2)/\pl^2$ & $0$ \\
\bottomrule
\end{tabular}
\end{center}

\subsection{Monotonicity}

\begin{proposition}[Reserve Monotonicity]\label{prop:reserve-monotone}
Both $x(p)$ and $y(p)$ are strictly increasing in $p$ on $[\pl, \pu]$:
\[
  x'(p) = \frac{2}{\pl^2} > 0, \qquad
  y'(p) = \frac{2p}{\pl^2} > 0 \quad\text{for } p > 0
\]
\end{proposition}

\begin{remark}
Both reserves increasing with $p$ is the correct direction under Framing~C
(\cref{sec:architecture}): as concentration rises, more virtual exposure
is absorbed ($x$ rises) and more numeraire accumulates ($y$ rises) because
the underwriter's collateral drains into the payoff liability. The LP value
$V(p) = y - px$ decreases, which is the underwriter bearing losses.
\end{remark}

\subsection{Collateral Accounting}

The underwriter deposits USDC (Token~Y). The virtual Token~X reserve is
an internal accounting variable with no external token.

\begin{definition}[Collateral Value]\label{def:collateral-value}
At price $p$, the underwriter's real (USDC) collateral per unit liquidity equals
the LP value:
\begin{equation}\label{eq:collateral-value}
  \text{collateral}(p) = V(p) = y(p) - p \cdot x(p) = \frac{\pu^2 - p^2}{\pl^2}
\end{equation}
At $p = \pl$: full collateral $(\pu^2 - \pl^2)/\pl^2$.
At $p = \pu$: zero (fully drained). The underwriter cannot lose more than
their initial deposit---the CLAMM tick cap $\pu$ enforces bounded loss.
\end{definition}

\subsection{Numerical Verification}

For $\DeltaStar = 0.09$, $\pl = 0.0989$, $\pu = 1.0$ ($\DeltaPlus_{\max} = 0.5$):

\begin{center}
\begin{tabular}{@{}lcccc@{}}
\toprule
$p$ & $\DeltaPlus$ & $x$ & $y$ & $V$ (USDC/unit $L$) \\
\midrule
$0.0989$ ($= \pl$) & $0.09$ & $20.22$ & $103.2$ & $101.2$ \\
$0.1758$ & $0.149$ & $35.95$ & $105.4$ & $99.04$ \\
$0.5$ & $0.333$ & $102.2$ & $127.6$ & $76.5$ \\
$1.0$ ($= \pu$) & $0.5$ & $204.5$ & $204.5$ & $0$ \\
\bottomrule
\end{tabular}
\end{center}

Payoff check at $p = 0.1758$: $N = V(\pl) - V(p) = 101.2 - 99.04 = 2.16$.
Formula: $(0.1758/0.0989)^2 - 1 = 3.16 - 1 = 2.16$ \checkmark.

% borrow-rate.tex — Utilization-based borrow rate (JumpRate) + separate funding
% Kontrol: prove_ins_borrowRate*, prove_ins_utilizationBound*

\section{Borrow Rate}\label{sec:borrow-rate}

The PLP (borrower) pays two separate rates to the underwriter (lender):
a unidirectional borrow rate (always PLP $\to$ underwriter) and a bidirectional
funding rate (\cref{sec:funding-rate}). These are economically distinct mechanisms
with separate accumulators, following the Primer on Perpetuals
\cite{angeris2022perpetuals}.

\subsection{Borrow Rate: JumpRate Curve}

\begin{definition}[Borrow Rate]\label{def:borrow-rate}
\begin{equation}\label{eq:borrow-rate}
  \rborrow(\Ua) = \premfactor + \mlow \cdot \min(\Ua, \kink)
                 + \mhigh \cdot \max(0,\, \Ua - \kink)
\end{equation}
where:
\begin{itemize}
  \item $\premfactor > 0$: base premium (floor), calibrated from econometric WTP;
  \item $\Ua = L_{\text{borrowed}} / L_{\text{total}} \in [0, 1]$: utilization ratio;
  \item $\kink \in (0, 1)$: target utilization (e.g., $\kink = 0.8$);
  \item $\mlow > 0$: slope below kink (gentle increase);
  \item $\mhigh \gg \mlow$: slope above kink (jump multiplier, strong deterrent).
\end{itemize}
\end{definition}

\begin{proposition}[Borrow Rate Properties]\label{prop:borrow-properties}
\begin{enumerate}[label=(\roman*)]
  \item $\rborrow \geq \premfactor > 0$ for all $\Ua \geq 0$ (underwriters always earn at least $\premfactor$).
  \item $\rborrow$ is continuous and piecewise-linear in $\Ua$.
  \item $\rborrow$ is non-decreasing in $\Ua$ (higher utilization $\to$ higher borrow cost).
  \item At $\Ua = 1$: $\rborrow = \premfactor + \mlow \kink + \mhigh(1 - \kink)$
    (very expensive, incentivizing new underwriter deposits or PLP exit).
\end{enumerate}
\end{proposition}

\subsection{Borrow Rate Accrual}

\begin{definition}[Borrow Fee Accumulator]\label{def:borrow-accumulator}
At each hook callback (afterSwap, afterAddLiquidity, afterRemoveLiquidity),
the borrow fee growth accumulator updates:
\begin{equation}\label{eq:borrow-accumulator}
  g^{\text{borrow}} \mathrel{+}= \rborrow(\Ua) \cdot \Delta t
\end{equation}
where $\Delta t$ is the time elapsed since the last callback (in blocks or seconds).
This accumulator is global (per insurance pool) and distributed to underwriters
pro-rata by liquidity share, using the Synthetix staking pattern
\cite{batog2018scalable}.
\end{definition}

\subsection{Separate Accounting Principle}

\begin{remark}[Two Accumulators, Two Purposes]
The borrow rate and funding rate (\cref{sec:funding-rate}) use separate
accumulators:
\begin{itemize}
  \item $g^{\text{borrow}}$: always positive, PLP $\to$ underwriter.
    Compensates capital lockup and tail risk.
  \item $g^{\text{funding}}$: bidirectional. Corrects mark-index mispricing.
    When positive, PLP pays; when negative, underwriter pays.
\end{itemize}
The Primer on Perpetuals \cite[Section~3]{angeris2022perpetuals} establishes that
conflating these mechanisms breaks funding convergence: a funding term that
cannot go negative (because it's floored by the borrow rate) cannot fully
correct underpricing.
\end{remark}

% funding-rate.tex — Bidirectional funding rate with jump adjustment
% Kontrol: prove_cfmm_perpetual_fundingConvergence, prove_cfmm_perpetual_jumpSafety

\section{Funding Rate}\label{sec:funding-rate}

This instrument is perpetual: it has no fixed expiry but settles at liquidity events.
The funding rate maintains mark-index convergence, following the framework of
\cite{angeris2022perpetuals}. It is separate from the borrow rate
(\cref{sec:borrow-rate}) and bidirectional.

\subsection{Mark and Index Prices}

\begin{definition}[Index Price]\label{def:index-price}
The index price is derived from the FCI hook's co-primary state:
\begin{equation}\label{eq:index-price}
  \pindex = \frac{\DeltaPlus}{1 - \DeltaPlus}
\end{equation}
where $\DeltaPlus = \max(0,\, \AT - \ATnull)$ is the realized concentration deviation
from the underlying pool. Computed on the fly from stored $\AT$, $\ThetaSum$, $\PosCount$.
\end{definition}

\begin{definition}[Mark Price]\label{def:mark-price}
The mark price is the CFMM's implied price from virtual reserves:
\begin{equation}\label{eq:mark-price}
  \pmark = \frac{\pl^2}{2}\sqrt{\frac{u}{\Lact}}
\end{equation}
where $u$ is the current virtual concentration reserve and $\Lact$ is
the active liquidity in the current tick range.
\end{definition}

\subsection{Discrete Funding at Hook Callbacks}

Funding is exchanged at each hook callback (\texttt{afterSwap},
\texttt{afterAddLiquidity}, \texttt{afterRemoveLiquidity}), matching
the FCI oracle update frequency. No keeper or per-block updates required.

\begin{definition}[Funding Payment]\label{def:funding-payment}
At callback $n$, the funding rate is:
\begin{equation}\label{eq:funding-rate}
  \rfunding_n = \alpha \cdot \frac{\pmark_n - \pindex_n}{\pindex_n + 1}
\end{equation}
where $\alpha > 0$ is the sensitivity parameter (set at pool initialization).
The denominator $(\pindex + 1)$ normalizes by index level, preventing
extreme rates when $\pindex$ is small.
\end{definition}

\begin{remark}[Bidirectionality]
\begin{itemize}
  \item $\pmark > \pindex$ (protection overpriced): $\rfunding > 0$,
    PLPs pay underwriters. Incentivizes arbitrageurs to sell the perpetual,
    pushing $\pmark$ down.
  \item $\pmark < \pindex$ (protection underpriced): $\rfunding < 0$,
    underwriters pay PLPs. Incentivizes buying the perpetual,
    pushing $\pmark$ up.
\end{itemize}
This bidirectionality is essential for convergence and is why funding must
be in a separate accumulator from the unidirectional borrow rate
(\cref{sec:borrow-rate}).
\end{remark}

\subsection{Jump-Adjusted Replication}

The backtest (ETH/USDC 30bps, 41 days) confirms that $\DeltaPlus$ exhibits
rare, large, discontinuous jumps: $0 \to 0.158$ in a single day.
Between callbacks, a single large swap or JIT event can change the FCI
dramatically. The funding rate includes a jump premium.

\begin{definition}[Jump Premium]\label{def:jump-premium}
When $|\Delta\pindex| > \delta_{\text{jump}}$ at callback $n$
(where $\delta_{\text{jump}}$ is the jump threshold), the funding
rate receives an additive adjustment:
\begin{equation}\label{eq:jump-funding}
  \rfunding_n^{\text{adj}} = \rfunding_n + \lambda \cdot |\Delta\pindex|
\end{equation}
where $\lambda > 0$ is the jump intensity parameter. This compensates
underwriters for unhedgeable gap risk between callbacks.
\end{definition}

\begin{theorem}[Jump Safety]\label{thm:jump-safety}
When $\DeltaPlus$ jumps discretely from $\delta$ to $\delta'$ at a liquidity event
(due to abrupt changes in $\PosCount$ or $\ThetaSum$), the protection value change
is bounded:
\begin{equation}\label{eq:jump-bound}
  |\Delta\text{value}| \leq |V(\delta') - V(\delta)|
\end{equation}
where $V$ is the LP payoff function (\cref{eq:lp-payoff}).
\end{theorem}

\begin{proof}
The LP payoff $V(p)$ is Lipschitz continuous on $[\pl, \pu]$ with
$|V'(p)| = 2p/\pl^2 \leq 2\pu/\pl^2$. The value change at a jump
equals $V(p(\delta')) - V(p(\delta))$, which is bounded by the payoff
difference because $V$ is monotone and no intermediate states are
observable between callbacks. The jump-adjusted funding compensates
for this discontinuity following
\cite[Section~3]{angeris2022perpetuals}. \qedhere
\end{proof}

\subsection{Funding Accumulator}

\begin{definition}[Funding Growth]\label{def:funding-accumulator}
The global funding growth accumulator:
\begin{equation}\label{eq:funding-accumulator}
  g^{\text{funding}} \mathrel{+}= \rfunding_n^{\text{adj}} \cdot \Delta t_n
\end{equation}
This accumulator can be positive or negative. Per-position funding owed
is computed via the Synthetix difference pattern:
$\text{funding}_i = L_i \cdot (g^{\text{funding}} - g^{\text{funding}}_{\text{baseline},i})$.
\end{definition}

\subsection{Convergence Property}

\begin{proposition}[Mark-Index Convergence]\label{prop:convergence}
Under the funding mechanism, $|\pmark - \pindex|$ decreases
after each callback (in expectation), provided $\alpha > 0$.
\end{proposition}

\begin{proof}[Proof sketch]
When $\pmark > \pindex$, PLPs pay positive funding, reducing demand
for the perpetual and pushing $\pmark$ down. When $\pmark < \pindex$,
underwriters pay, increasing demand and pushing $\pmark$ up.
The sensitivity $\alpha$ controls convergence speed.
A formal convergence proof under specific assumptions on arrival rates
is deferred to the Kontrol verification phase. \qedhere
\end{proof}

% initialization.tex — Insurance CFMM initialization and registration
% Kontrol: prove_ins_init*, prove_ins_registration*

\section{Initialization}\label{sec:initialization}

\subsection{Pool Initialization}

\begin{definition}[Insurance Pool Parameters]\label{def:init-params}
An insurance CFMM is initialized per underlying V4 pool with:
\begin{enumerate}[label=(\roman*)]
  \item $\texttt{fciHook}$: address of the deployed FeeConcentrationIndex hook
    (oracle for $\pindex$);
  \item $\texttt{poolManager}$: Uniswap V4 PoolManager address;
  \item $\premfactor \in (0, 1)$: base premium factor (borrow rate floor);
  \item $\alpha > 0$: funding rate sensitivity;
  \item $\lambda > 0$: jump intensity parameter;
  \item $\kink \in (0, 1)$: utilization target for JumpRate curve;
  \item $\mlow, \mhigh > 0$: JumpRate slope parameters ($\mhigh \gg \mlow$);
  \item $\texttt{tickSpacing}$: tick spacing for the insurance CFMM tick array.
\end{enumerate}
These parameters are immutable after deployment.
\end{definition}

\subsection{First Underwriter Deposit}

\begin{definition}[Initialization from Oracle]\label{def:init-oracle}
The first underwriter deposits $D$ USDC into tick range $[\pl, \pu]$.
The FCI hook provides current $\DeltaPlus$, yielding $p_0 = \DeltaPlus/(1-\DeltaPlus)$.

\textbf{Case 1:} $p_0 \leq \pl$ (no active concentration---typical at launch):
\begin{align}
  L_0 &= \frac{D \cdot \pl^2}{\pu^2 + \pl^2} \label{eq:init-L-cold} \\
  u_0 &= L_0 \cdot \frac{4\pl^2}{\pl^4} = \frac{4L_0}{\pl^2} \\
  y_0 &= D
\end{align}
Spot price initializes at $\pl$ (the strike boundary).

\textbf{Case 2:} $\pl < p_0 < \pu$ (concentration already active):
\begin{align}
  L_0 &= \frac{D \cdot \pl^2}{\pu^2 + p_0^2} \label{eq:init-L-warm} \\
  u_0 &= L_0 \cdot \frac{4p_0^2}{\pl^4} \\
  y_0 &= D
\end{align}
The underwriter enters mid-range with immediate concentration liability.

\textbf{Case 3:} $p_0 \geq \pu$ (concentration above cap):
Deposit reverts. The underwriter would face instant max-loss.
\end{definition}

\begin{proposition}[Invariant at Initialization]\label{prop:init-invariant}
In all valid cases: $\psi(u_0, y_0) = L_0 \cdot \pu^2/\pl^2$.
\end{proposition}

\begin{proof}
For Case~1: $\psi = D - (\pl^2/4)(4L_0/\pl^2) = D - L_0
= D - D\pl^2/(\pu^2 + \pl^2) = D \cdot \pu^2/(\pu^2 + \pl^2)
= L_0 \cdot (\pu^2 + \pl^2)/\pl^2 \cdot \pu^2/(\pu^2 + \pl^2)
= L_0 \cdot \pu^2/\pl^2$. \qedhere
\end{proof}

\subsection{PLP Registration (Borrowing)}

\begin{definition}[PLP Borrows LP Share]\label{def:plp-borrow}
A PLP holding a V4 LP position borrows insurance CFMM LP shares:
\begin{enumerate}[label=(\roman*)]
  \item PLP calls \texttt{borrow(poolKey, v4PositionId, borrowAmount)}.
  \item The hook verifies the PLP holds a valid V4 position with active fee stream.
  \item LP shares are transferred from the insurance CFMM to the PLP
    (reducing $L_{\text{available}}$, increasing $L_{\text{borrowed}}$).
  \item The PLP's $\max_t p(t)$ tracker is initialized at current $\pindex$.
  \item Utilization $\Ua$ increases, raising $\rborrow$ for all borrowers.
\end{enumerate}
\end{definition}

\subsection{PLP Exit (Returning LP Share)}

\begin{definition}[Exit Settlement]\label{def:plp-exit}
At \texttt{afterRemoveLiquidity} (or voluntary deregistration):
\begin{enumerate}[label=(\roman*)]
  \item Compute $\text{premium}_i = \premfactor \times \text{lifetime\_fees}_i
    + \text{accrued\_borrow\_rate}_i + \text{accrued\_funding}_i$.
  \item Compute payout: $\text{payout}_i = \text{premium}_i \times
    [(\max_t p(t) / \pl)^2 - 1]^{+}$.
  \item The payout is capped at $\text{premium}_i \times [(\pu/\pl)^2 - 1]$
    (tick range cap).
  \item LP shares are returned to the insurance CFMM. Utilization $\Ua$ decreases.
  \item Net settlement: PLP receives $\max(\text{payout} - \text{premium}_i, 0)$
    from the CFMM reserves. If payout $<$ premium, the difference stays in reserves.
\end{enumerate}
\end{definition}

\begin{remark}[Max-Price Tracking]
Each PLP position stores $p_{\max,i} = \max_{t \in [\text{mint}, \text{now}]} \pindex(t)$.
Updated at every hook callback: $p_{\max,i} \gets \max(p_{\max,i}, \pindex)$.
This is one SSTORE per position per callback (only if $\pindex > p_{\max,i}$),
amortized by checking the condition before writing.
\end{remark}

% invariants.tex — Complete invariant set for the fee concentration insurance CFMM
% Kontrol: prove_cfmm_*, prove_fci_*, prove_ins_*

\section{Invariants}\label{sec:invariants}

This section contains all invariants for the fee concentration insurance system.
The invariant set is organized into five groups:
\begin{enumerate}
  \item CFMM-01 through CFMM-19: Base CFMM invariants (trading set, reserves, fees, ticks).
  \item CFMM-20 through CFMM-25: Replication invariants (payoff, reserves, canonical form).
  \item CFMM-26 through CFMM-28: Perpetual invariants (funding, jump safety).
  \item FCI-01 through FCI-13: Fee concentration index invariants (state, price mapping, perpetual funding).
  \item INS-01 through INS-10: Insurance CFMM invariants (borrow, settlement, utilization).
\end{enumerate}

% ══════════════════════════════════════════════════════════════════════════════
\subsection{Base CFMM Invariants (CFMM-01 through CFMM-19)}\label{sec:base-cfmm}
% ══════════════════════════════════════════════════════════════════════════════

\begin{invariant}{CFMM-01}{Trading set convexity}
% Kontrol: prove_cfmm_tradingSetConvexity
\[
  S = \{(u, y) : \psi(u, y) \geq L \cdot C_1\} \text{ is convex}
\]
The trading set is convex. Since $\psi(u,y) = y - (\pl^2/4)u$ is linear (hence concave),
the superlevel set is convex (Geometry paper \cite{angeris2023geometry}, Theorem~1).
\end{invariant}

\begin{invariant}{CFMM-02}{Reserve non-negativity}
% Kontrol: prove_cfmm_reserveNonNeg
\[
  u \geq 0 \quad\text{and}\quad y \geq 0
\]
Both virtual concentration reserve $u$ and USDC reserve $y$ are non-negative at all times.
\end{invariant}

\begin{invariant}{CFMM-03}{Invariant preservation on swap}
% Kontrol: prove_cfmm_invariantSwap
\[
  \{\psi(u, y) = L \cdot C_1\}
  \;\xrightarrow{\text{swap}(\Delta u)}\;
  \{\psi(u + \Delta u, y - \Delta y) = L \cdot C_1\}
\]
Every swap preserves the trading function invariant. From $\Delta y = (\pl^2/4)\Delta u$.
\end{invariant}

\begin{invariant}{CFMM-04}{Invariant preservation on mint}
% Kontrol: prove_cfmm_invariantMint
\[
  \{\psi(u, y) = L \cdot C_1\}
  \;\xrightarrow{\text{mint}(\Delta L)}\;
  \{\psi(u + \Delta u, y + \Delta y) = (L + \Delta L) \cdot C_1\}
\]
Minting new liquidity preserves the invariant at the new liquidity level
(\cref{prop:mint-invariant}).
\end{invariant}

\begin{invariant}{CFMM-05}{Invariant preservation on burn}
% Kontrol: prove_cfmm_invariantBurn
\[
  \{\psi(u, y) = L \cdot C_1\}
  \;\xrightarrow{\text{burn}(\Delta L)}\;
  \{\psi(u - \Delta u, y - \Delta y) = (L - \Delta L) \cdot C_1\}
\]
Burning liquidity preserves the invariant at the reduced liquidity level.
\end{invariant}

\begin{invariant}{CFMM-06}{Spot price consistency}
% Kontrol: prove_cfmm_spotPriceConsistency
\[
  p = \frac{\pl^2}{2}\sqrt{\frac{u}{\Lact}}
  = -\frac{\partial\psi/\partial u}{\partial\psi/\partial y} \cdot \frac{du}{dx}
\]
The spot price derived from reserves matches the marginal rate of substitution
(\cref{def:spot-price}).
\end{invariant}

\begin{invariant}{CFMM-07}{Price within tick range}
% Kontrol: prove_cfmm_priceInRange
\[
  \pl \leq p \leq \pu \quad\text{when } \Lact > 0
\]
The spot price stays within the active tick range $[\pl, \pu]$
when there is active liquidity.
\end{invariant}

\begin{invariant}{CFMM-08}{Liquidity additivity at tick crossing}
% Kontrol: prove_cfmm_tickCrossing
\[
  \Lact' = \Lact + \Lnet
\]
At each tick crossing, active liquidity updates by adding the signed net liquidity
of the crossed tick. This is O(1) because $\psi$ is 1-homogeneous in $L$
(\cref{thm:linearity}).
\end{invariant}

\begin{invariant}{CFMM-09}{Liquidity non-negativity}
% Kontrol: prove_cfmm_liquidityNonNeg
\[
  L \geq 0 \quad\text{and}\quad \Lact \geq 0
\]
Total and active liquidity are non-negative.
\end{invariant}

\begin{invariant}{CFMM-10}{No-arbitrage (monotone reserves)}
% Kontrol: prove_cfmm_noArbitrage
\[
  p' > p \;\longrightarrow\; u(p') > u(p) \;\land\; y(p') > y(p)
\]
Both reserves are strictly increasing in price (\cref{prop:reserve-monotone}).
Under Framing~C, rising concentration absorbs more virtual exposure and
accumulates more USDC.
\end{invariant}

\begin{invariant}{CFMM-11}{Fee accumulator monotonicity}
% Kontrol: prove_cfmm_feeAccMonotone
\[
  g^{\text{borrow}}_{n+1} \geq g^{\text{borrow}}_n
\]
The borrow fee growth accumulator is non-decreasing (borrow rate $\geq \premfactor > 0$).
\end{invariant}

\begin{invariant}{CFMM-12}{Fee accumulator per-position}
% Kontrol: prove_cfmm_feeAccPerPosition
\[
  \text{fees}_i = L_i \cdot (g^{\text{borrow}} - g^{\text{borrow}}_{\text{baseline},i})
\]
Per-position borrow fees use the Synthetix difference pattern \cite{batog2018scalable}.
\end{invariant}

\begin{invariant}{CFMM-13}{LP value bounded}
% Kontrol: prove_cfmm_lpValueBounded
\[
  0 \leq V(p) = \frac{\pu^2 - p^2}{\pl^2} \leq \frac{\pu^2 - \pl^2}{\pl^2}
\]
LP value is bounded between 0 (at $p = \pu$, full loss) and the initial
collateral (at $p = \pl$, no loss).
\end{invariant}

\begin{invariant}{CFMM-14}{Swap output positive}
% Kontrol: prove_cfmm_swapOutputPositive
\[
  \Delta x > 0 \;\longrightarrow\; \Delta y > 0
  \quad\text{and}\quad
  \Delta x < 0 \;\longrightarrow\; \Delta y < 0
\]
Swaps in either direction produce correctly-signed outputs.
\end{invariant}

\begin{invariant}{CFMM-15}{Slippage non-negative}
% Kontrol: prove_cfmm_slippageNonNeg
\[
  \frac{\Delta y}{\Delta x} = \frac{\pl^2}{4}(2x + \Delta x) \geq \frac{\pl^2}{4} \cdot 2x = p
\]
The effective price is always at least the spot price (natural slippage from
the $u = x^2$ transform). Equality only for infinitesimal trades.
\end{invariant}

\begin{invariant}{CFMM-16}{Coordinate transform bijection}
% Kontrol: prove_cfmm_coordTransformBijection
\[
  u = x^2 \;\longleftrightarrow\; x = \sqrt{u} \quad\text{for } x \geq 0
\]
The $u \leftrightarrow x$ mapping is a bijection on the non-negative reals.
One \texttt{FixedPointMathLib.sqrt} call per swap.
\end{invariant}

\begin{invariant}{CFMM-17}{1-homogeneity}
% Kontrol: prove_cfmm_oneHomogeneity
\[
  \psi(\lambda u, \lambda y) = \lambda \psi(u, y) \quad\forall\, \lambda > 0
\]
The trading function is 1-homogeneous, enabling O(1) tick-based liquidity
aggregation (\cref{thm:linearity}).
\end{invariant}

\begin{invariant}{CFMM-18}{Concavity}
% Kontrol: prove_cfmm_concavity
\[
  \nabla^2 \psi =
  \begin{pmatrix} 0 & 0 \\ 0 & 0 \end{pmatrix}
\]
$\psi(u,y) = y - (\pl^2/4)u$ is linear, hence (weakly) concave.
The Hessian is zero. The trading set is convex.
\end{invariant}

\begin{invariant}{CFMM-19}{Single-sided deposit consistency}
% Kontrol: prove_cfmm_singleSidedDeposit
\[
  \Delta L = \frac{\Delta y \cdot \pl^2}{\pu^2 + p^2}
  \quad\text{and}\quad
  \Delta u = \Delta L \cdot \frac{4p^2}{\pl^4}
\]
Single-sided USDC deposits correctly compute $\Delta L$ and virtual $\Delta u$
(\cref{def:single-sided}). Only USDC enters the contract.
\end{invariant}

% ══════════════════════════════════════════════════════════════════════════════
\subsection{Replication Invariants (CFMM-20 through CFMM-25)}\label{sec:replication}
% ══════════════════════════════════════════════════════════════════════════════

Source: Angeris et al.\ \cite{angeris2021replicating,angeris2023geometry}.

\begin{invariant}{CFMM-20}{Replication accuracy}
% Kontrol: prove_cfmm_replicationAccuracy
\[
  \{V(p)\}
  \;\xrightarrow{\text{evolve}(p \to p')}\;
  \{|\text{LP\_value}(p') - V(p')| \leq \varepsilon\}
\]
The CFMM LP position replicates the target payoff $V(p) = (\pu^2 - p^2)/\pl^2$
within error $\varepsilon$ (determined by discretization and fee leakage).
\end{invariant}

\begin{invariant}{CFMM-21}{LP payoff concavity}
% Kontrol: prove_cfmm_payoffConcavity
\[
  V''(p) = -\frac{2}{\pl^2} < 0 \quad\forall\, p \in [\pl, \pu]
\]
The LP payoff is strictly concave on the active domain.
Required for Angeris replication theorem \cite{angeris2021replicating}.
\end{invariant}

\begin{invariant}{CFMM-22}{Reserve-payoff consistency}
% Kontrol: prove_cfmm_reservePayoffConsistency
\[
  x(p) = -V'(p) = \frac{2p}{\pl^2}
  \quad\text{and}\quad
  y(p) = V(p) + p \cdot x(p) = \frac{\pu^2 + p^2}{\pl^2}
\]
Reserves are the derivatives and Legendre transform of the payoff
(\cref{def:reserves}).
\end{invariant}

\begin{invariant}{CFMM-23}{Boundary reserves}
% Kontrol: prove_cfmm_boundaryReserves
\begin{align*}
  p = \pl &:\quad x = 2/\pl,\; y = (\pu^2 + \pl^2)/\pl^2,\; V = (\pu^2 - \pl^2)/\pl^2 \\
  p = \pu &:\quad x = 2\pu/\pl^2,\; y = 2\pu^2/\pl^2,\; V = 0
\end{align*}
At the lower boundary, full collateral remains. At the upper boundary, collateral
is fully drained (max loss for underwriter, max payout for PLP).
\end{invariant}

\begin{invariant}{CFMM-24}{Price-reserve consistency}
% Kontrol: prove_cfmm_priceReserveConsistency
\[
  p = \frac{\pl^2}{2} \cdot \frac{x}{L}
  \quad\Longleftrightarrow\quad
  x(p) = \frac{2pL}{\pl^2}
\]
The price derived from the trading function marginal rate matches the
reserve curve parameterization (\cref{def:spot-price}).
\end{invariant}

\begin{invariant}{CFMM-25}{Canonical form}
% Kontrol: prove_cfmm_canonicalForm
\[
  \psi(u, y) = y - \frac{\pl^2}{4}\, u
\]
is nondecreasing in $y$, nonincreasing in $u$, concave, and 1-homogeneous.
This is the unique canonical trading function for this CFMM
(Geometry paper \cite{angeris2023geometry}, Theorem~1).
\end{invariant}

% ══════════════════════════════════════════════════════════════════════════════
\subsection{Perpetual Invariants (CFMM-26 through CFMM-28)}\label{sec:perpetual-invariants}
% ══════════════════════════════════════════════════════════════════════════════

Source: Angeris et al.\ \cite{angeris2022perpetuals}.

\begin{invariant}{CFMM-26}{Funding rate convergence}
% Kontrol: prove_cfmm_perpetual_fundingConvergence
\[
  \{|\pmark - \pindex| > \delta\}
  \;\xrightarrow{\text{funding}(\Delta t)}\;
  \{|\pmark' - \pindex'| < |\pmark - \pindex|\}
\]
The bidirectional funding mechanism drives mark-index convergence
(\cref{prop:convergence}). Sensitivity $\alpha > 0$ controls speed.
\end{invariant}

\begin{invariant}{CFMM-27}{Funding rate bounds}
% Kontrol: prove_cfmm_perpetual_fundingBounds
\[
  |\rfunding_n| \leq \alpha + \lambda \cdot |\Delta\pindex|_{\max}
\]
The funding rate (including jump premium) is bounded. The bound depends on
the sensitivity $\alpha$, jump intensity $\lambda$, and maximum observable
$\pindex$ change between callbacks.
\end{invariant}

\begin{invariant}{CFMM-28}{Jump-adjusted value bound}
% Kontrol: prove_cfmm_perpetual_jumpSafety
\[
  \{\DeltaPlus \text{ jumps from } \delta \text{ to } \delta'\}
  \;\longrightarrow\;
  \{|\Delta\text{value}| \leq |V(\delta') - V(\delta)|\}
\]
Discrete jumps in $\DeltaPlus$ at liquidity events produce value changes
bounded by the payoff difference (\cref{thm:jump-safety}).
\end{invariant}

% ══════════════════════════════════════════════════════════════════════════════
\subsection{Fee Concentration State Invariants (FCI-01 through FCI-08)}\label{sec:fci-state}
% ══════════════════════════════════════════════════════════════════════════════

\begin{invariant}{FCI-01}{Index boundedness}
% Kontrol: prove_fci_indexBoundedness
\[
  0 \leq \AT \leq 1
\]
The fee concentration index is bounded in $[0,1]$ at all times.
$\AT = 0$ when no positions have been removed. $\AT = 1$ is the theoretical
maximum (single JIT monopoly).
\end{invariant}

\begin{invariant}{FCI-02}{Theta sum non-negativity}
% Kontrol: prove_fci_thetaSumNonNeg
\[
  \ThetaSum \geq 0
\]
The aggregate turnover rate $\ThetaSum = \sum_k \thetapos_k$ is non-negative
since each $\thetapos_k = 1/\blocklife_k > 0$.
\end{invariant}

\begin{invariant}{FCI-03}{Position count non-negativity}
% Kontrol: prove_fci_posCountNonNeg
\[
  \PosCount \geq 0
\]
The active position count is a non-negative integer.
\end{invariant}

\begin{invariant}{FCI-04}{Null lower bound}
% Kontrol: prove_fci_nullLowerBound
\[
  \{\PosCount > 0\} \;\longrightarrow\; \{\AT \geq \ATnull\}
\]
The real fee concentration index is always at least as large as the Ma--Crapis
equal-share null (\cref{prop:null-lower-bound}). Equivalently, $\DeltaPlus \geq 0$.
\end{invariant}

\begin{invariant}{FCI-05}{Deviation non-negativity}
% Kontrol: prove_fci_deviationNonNeg
\[
  \DeltaPlus = \max(0,\, \AT - \ATnull) \geq 0
\]
By construction ($\max$ with zero) and by \cref{prop:null-lower-bound}.
\end{invariant}

\begin{invariant}{FCI-06}{Deviation upper bound}
% Kontrol: prove_fci_deviationUpperBound
\[
  \DeltaPlus < 1
\]
Since $\AT \leq 1$ and $\ATnull \geq 0$, the deviation satisfies $\DeltaPlus \leq 1$.
Strict inequality holds when $\PosCount \geq 1$ (at least one position implies $\ATnull > 0$).
This ensures the price mapping $p = \DeltaPlus/(1-\DeltaPlus)$ is well-defined.
\end{invariant}

\begin{invariant}{FCI-07}{Co-primary consistency}
% Kontrol: prove_fci_coPrimaryConsistency
\[
  \ATnull = \sqrt{\frac{\ThetaSum}{\PosCount^2}}
  \quad\text{and}\quad
  \DeltaPlus = \max(0,\, \AT - \ATnull)
\]
The derived quantities $\ATnull$ and $\DeltaPlus$ are deterministic functions of
the three stored state variables $(\AT, \ThetaSum, \PosCount)$.
No additional state is needed.
\end{invariant}

\begin{invariant}{FCI-08}{Liquidity event atomicity}
% Kontrol: prove_fci_liquidityEventAtomicity
\[
  \{\texttt{afterRemoveLiquidity}\}
  \;\longrightarrow\;
  \{\AT \text{ updated} \;\land\; \ThetaSum \text{ updated} \;\land\; \PosCount \text{ updated}\}
\]
All three co-primary state variables update atomically at each liquidity event.
No intermediate state is observable where some variables are updated and others are not.
\end{invariant}

% ══════════════════════════════════════════════════════════════════════════════
\subsection{Price Mapping Invariants (FCI-09 through FCI-11)}
% ══════════════════════════════════════════════════════════════════════════════

\begin{invariant}{FCI-09}{Price non-negativity}
% Kontrol: prove_fci_priceNonNeg
\[
  p = \frac{\DeltaPlus}{1 - \DeltaPlus} \geq 0
\]
Since $\DeltaPlus \geq 0$ (FCI-05) and $1 - \DeltaPlus > 0$ (FCI-06).
\end{invariant}

\begin{invariant}{FCI-10}{Price monotonicity}
% Kontrol: prove_fci_priceMonotonicity
\[
  \DeltaPlus_1 < \DeltaPlus_2
  \;\longrightarrow\;
  p(\DeltaPlus_1) < p(\DeltaPlus_2)
\]
The price mapping is strictly increasing (\cref{prop:price-monotone}).
Higher concentration deviation $\to$ higher insurance price.
\end{invariant}

\begin{invariant}{FCI-11}{Price-deviation invertibility}
% Kontrol: prove_fci_priceInvertibility
\[
  \DeltaPlus = \frac{p}{1 + p}
  \quad\text{is the unique inverse of}\quad
  p = \frac{\DeltaPlus}{1 - \DeltaPlus}
\]
The mapping is a bijection from $[0,1) \to [0,\infty)$.
\end{invariant}

% ══════════════════════════════════════════════════════════════════════════════
\subsection{FCI Perpetual Funding Invariants (FCI-12 through FCI-13)}
% ══════════════════════════════════════════════════════════════════════════════

Source: Angeris et al.\ \cite{angeris2022perpetuals}.

\begin{invariant}{FCI-12}{Funding rate convergence (FCI)}
% Kontrol: prove_fci_fundingConvergence
\[
  \{|\pmark - \pindex| > \delta\}
  \;\longrightarrow\; \text{funding}(\text{liquidityEvent})
  \;\longrightarrow\; \{|\pmark' - \pindex'| < |\pmark - \pindex|\}
\]
The funding mechanism drives mark-index convergence (\cref{prop:convergence}).
\end{invariant}

\begin{invariant}{FCI-13}{Jump-adjusted value bound (FCI)}
% Kontrol: prove_fci_jumpBound
\[
  \{\DeltaPlus \text{ jumps from } \delta \text{ to } \delta'\}
  \;\longrightarrow\;
  \{|\Delta\text{value}| \leq |V(\delta') - V(\delta)|\}
\]
Discrete jumps in $\DeltaPlus$ at liquidity events do not cause value changes
exceeding the payoff difference (\cref{thm:jump-safety}).
\end{invariant}

% ══════════════════════════════════════════════════════════════════════════════
\subsection{Insurance CFMM Invariants (INS-01 through INS-10)}\label{sec:insurance-cfmm}
% ══════════════════════════════════════════════════════════════════════════════

\begin{invariant}{INS-01}{Collateral solvency}
% Kontrol: prove_ins_collateralSolvency
\[
  y \geq L \cdot \frac{\pu^2 + p^2}{\pl^2} \quad\text{(USDC reserve covers all LP positions)}
\]
The USDC reserve is sufficient to cover the numeraire reserve requirement
at the current price. Underwriters cannot lose more than their deposit
(\cref{def:collateral-value}).
\end{invariant}

\begin{invariant}{INS-02}{Utilization bound}
% Kontrol: prove_ins_utilizationBound
\[
  0 \leq \Ua = \frac{L_{\text{borrowed}}}{L_{\text{total}}} \leq 1
\]
Utilization is bounded in $[0,1]$. $L_{\text{borrowed}} \leq L_{\text{total}}$
is enforced by reverting borrows that would exceed available liquidity.
\end{invariant}

\begin{invariant}{INS-03}{Borrow rate positivity}
% Kontrol: prove_ins_borrowRatePositivity
\[
  \rborrow(\Ua) \geq \premfactor > 0 \quad\forall\, \Ua \in [0,1]
\]
The borrow rate is always at least $\premfactor$ (the base premium floor).
Underwriters always earn a positive rate (\cref{prop:borrow-properties}).
\end{invariant}

\begin{invariant}{INS-04}{Borrow rate monotonicity}
% Kontrol: prove_ins_borrowRateMonotonicity
\[
  \Ua_1 < \Ua_2 \;\longrightarrow\; \rborrow(\Ua_1) \leq \rborrow(\Ua_2)
\]
The borrow rate is non-decreasing in utilization. Higher utilization
means higher cost, incentivizing new underwriter deposits or PLP exit
(\cref{prop:borrow-properties}).
\end{invariant}

\begin{invariant}{INS-05}{Separate accumulator accounting}
% Kontrol: prove_ins_separateAccumulators
\[
  g^{\text{borrow}} \geq 0
  \quad\text{(always positive)}
  \qquad
  g^{\text{funding}} \in \mathbb{R}
  \quad\text{(bidirectional)}
\]
The borrow and funding accumulators are maintained separately
(\cref{def:borrow-accumulator,def:funding-accumulator}).
The borrow accumulator is non-negative (PLP always pays underwriter).
The funding accumulator can be positive or negative (bidirectional).
Conflating them breaks funding convergence \cite{angeris2022perpetuals}.
\end{invariant}

\begin{invariant}{INS-06}{Max-price tracking}
% Kontrol: prove_ins_maxPriceTracking
\[
  p_{\max,i} = \max_{t \in [\text{mint}_i, \text{now}]} \pindex(t)
  \quad\text{and}\quad
  p_{\max,i}' \geq p_{\max,i}
\]
Each PLP position stores the running maximum of $\pindex$ since mint.
Updated at every hook callback when $\pindex > p_{\max,i}$.
This value is monotone non-decreasing.
\end{invariant}

\begin{invariant}{INS-07}{Payout formula}
% Kontrol: prove_ins_payoutFormula
\[
  \text{payout}_i = \text{premium}_i \times \left[\left(\frac{p_{\max,i}}{\pl}\right)^2 - 1\right]^{+}
\]
The PLP payout at exit follows the power-2 payoff. The plus operator
$[\cdot]^+$ ensures non-negative payouts. The premium is defined as
$\text{premium}_i = \premfactor \cdot \text{lifetime\_fees}_i
+ \text{accrued\_borrow}_i + \text{accrued\_funding}_i$.
\end{invariant}

\begin{invariant}{INS-08}{Payout cap}
% Kontrol: prove_ins_payoutCap
\[
  \text{payout}_i \leq \text{premium}_i \times \left[\left(\frac{\pu}{\pl}\right)^2 - 1\right]
\]
The payout is capped by the tick range. Since $p_{\max,i} \leq \pu$ within
the CLAMM range, the payout cannot exceed the cap. No global cap parameter
is needed---the per-position tick range enforces bounded loss for underwriters.
\end{invariant}

\begin{invariant}{INS-09}{Initialization invariant}
% Kontrol: prove_ins_initInvariant
\[
  \psi(u_0, y_0) = L_0 \cdot \frac{\pu^2}{\pl^2}
\]
At initialization (first deposit), the invariant holds with
$C_1 = \pu^2/\pl^2$ (\cref{prop:init-invariant}).
\end{invariant}

\begin{invariant}{INS-10}{Exit settlement conservation}
% Kontrol: prove_ins_exitSettlement
\[
  \{\texttt{exit}_i\}
  \;\longrightarrow\;
  \{L_{\text{borrowed}}' = L_{\text{borrowed}} - L_i\}
  \;\land\;
  \{\text{net}_i = \max(\text{payout}_i - \text{premium}_i, 0)\}
\]
At PLP exit: borrowed liquidity returns to the pool (utilization decreases),
and the PLP receives the net settlement (payout minus premium, floored at zero).
If payout $<$ premium, the excess stays in reserves as underwriter profit.
\end{invariant}


\bibliographystyle{alpha}
\bibliography{references}

\end{document}
